\documentclass[11pt]{article}
\usepackage[T1]{fontenc}
\usepackage{lmodern,amsmath,amssymb,amsthm,mathtools,mathrsfs,enumitem,booktabs}
\usepackage[margin=1in,headheight=15pt]{geometry}
\usepackage{microtype,xcolor,fancyhdr,listings}
\usepackage[colorlinks=true,linkcolor=teal!65!black,urlcolor=teal!65!black,citecolor=teal!65!black]{hyperref}
\hypersetup{pdftitle={iPRA with one explicit computation schema},
 pdfauthor={Prepared for Vladimir Reshetnikov},
 pdfsubject={Canonical successor-seeded syntax, bi-interpretability, and complete-checker accounting},
 pdfkeywords={iPRA, intuitionistic arithmetic, primitive recursion, axiomatization}}
\newtheorem{theorem}{Theorem}[section]
\newtheorem{lemma}[theorem]{Lemma}
\newtheorem{proposition}[theorem]{Proposition}
\newtheorem{corollary}[theorem]{Corollary}
\theoremstyle{definition}
\newtheorem{definition}[theorem]{Definition}
\newtheorem{remark}[theorem]{Remark}
\newcommand{\norm}[1]{\nu(#1)}
\newcommand{\subc}[2]{#1\langle #2\rangle}
\newcommand{\lit}{\mathrel{\equiv}}
\newcommand{\Rfree}{\mathcal R_{\mathrm{free}}}
\newcommand{\Cnorm}{\mathcal C_\nu}
\newcommand{\Knorm}{\mathcal K_\nu}
\newcommand{\ipra}{\mathsf{iPRA}}
\newcommand{\llbracket}{\mathopen{[\![}}
\newcommand{\rrbracket}{\mathclose{]\!]}}
\newcommand{\file}[1]{{\small\nolinkurl{#1}}}
\setcounter{tocdepth}{1}
\setlength{\emergencystretch}{2em}
\pagestyle{fancy}
\fancyhf{}
\fancyhead[L]{\small iPRA: canonical recursors}
\fancyhead[R]{\small November 2026}
\fancyfoot[C]{\thepage}
\lstset{basicstyle=\ttfamily\footnotesize,breaklines=true,columns=fullflexible,
 keepspaces=true,showstringspaces=false,frame=single,rulecolor=\color{gray!45}}
\title{\textbf{iPRA with one explicit computation schema}\\[0.4em]
\large Canonical successor-seeded syntax and complete-checker accounting}
\author{Prepared for Vladimir Reshetnikov}
\date{12 November 2026}
\begin{document}
\maketitle
\begin{abstract}
We give a presentation of intuitionistic primitive recursive arithmetic with
six inference patterns, two numerical constructors in addition to variables,
and one explicit computation schema. Successor is defined by a
successor-seeded recursor. Its zero computation is imposed as a canonical
spelling of terms, and normalized substitution is specified and charged.
We prove both translations to the existing arithmetic and both composites,
including the recursor body-extensionality lemma required by normalization.
A complete strict checker has 4,611 normalized bytes against 4,923 for a
matching seven-rule scalar control, a 6.34 percent reduction. A separate
516-byte refinement improves the existing compact guarded calculus and is
not assigned to the new axiomatization. We also prove two boundaries: a
scalar self-seeded language cannot define predecessor, and induction on
zero equations does not recover arbitrary equation induction, even with
coverage and successor injectivity. The latter uses a nonstandard
two-point-fiber countermodel. All new metatheorems are written proofs;
the archive includes executable tests, exact source controls, and
reproduction instructions, without claiming new proof-assistant verification
or an absolute minimum.
\end{abstract}
\clearpage
\tableofcontents
\newpage


\section{Result and comparison contract}
\label{sec:assessment}

The proposed presentation has six primitive inference patterns and the
numerical term constructors $0$ and $R$, in addition to variables. Successor
is a derived notation. There is one explicit computation schema. Its
ordinary realization is bi-interpretable with full first-order intuitionistic
primitive recursive arithmetic, with the entire numerical domain in both
directions and with identity graph $x=y$ for both numerical composites.

The reduction is precise but has a cost: the zero computation becomes a
canonical-syntax convention, and substitution must enforce that convention.
We count this normalization in the complete implementation. A reduction from
seven named proof rules to six is therefore not, on its own, a proof of a
smaller complete checker.

\begin{center}
\small
\begin{tabular}{lll}
\toprule
Feature & Current scalar presentation & Canonical $R$ presentation\\
\midrule
Numerical constructors & $0,S,I$ & $0,R$\\
Explicit computation schemas & Two & One\\
Primitive inference patterns & Seven & Six\\
Term equality in the checker & Literal syntax & Canonical syntax\\
Substitution & Capture avoiding & Capture avoiding, then normalize\\
Numerical interpretation domain & Whole & Whole\\
Numerical composite graph & $x=y$ & $x=y$\\
\bottomrule
\end{tabular}
\end{center}

The comparison is pinned to ProveIt commit
\texttt{8eb71c1e570e07ef45ead686628e77f18ef34798}, inspected on
12 November 2026, UTC. The current seven-rule scalar calculus is
\texttt{MPrf} in \texttt{NativeMinimal.lean}. The separate crossed calculus
\texttt{RPrf} is in \texttt{NativeReversal.lean}. Their equivalences with the
native scalar core are already formalized upstream. The present canonical
$R$ reduction has a written proof and executable implementation tests; it is
not a new Lean or Rocq formalization.

The repository's smallest compact profiles are a different comparison
category. The selected byte profile has 522 normalized source bytes, but its
bi-interpretation theorem concerns a guarded public language, with broader
auxiliary proof syntax. Its compact program does not itself validate that
public grammar. Our strict canonical-$R$ checker validates its full public
language. We therefore report two separate comparisons: matching strict
implementations of the mathematical calculi, and small implementation
changes to the existing compact calculus. No compact size is attributed to
the new six-rule mathematics.

The successor-seeded recursor, its predecessor construction, the scalar
arithmetic bridge, crossed transport, and witness-directed introduction are
inherited. The mathematical change developed here is canonicalization of the
successor-seeded syntax and the proof that it preserves substitution,
derivability, and both interpretation composites. A separate strengthening
of the scalar self-seeded obstruction is proved in
Section~\ref{sec:scalar-obstruction}.



\section{Scope and the exact mathematical change}

The comparison source is ProveIt commit
\texttt{8eb71c1e570e07ef45ead686628e77f18ef34798}.
The successor-seeded recursor and its arithmetic completeness are inherited from
\file{Research/successor_seed_completion/} and
\file{Research/scalar_successor/}.
The crossed-transport argument is the retraction criterion of
\file{Research/reversal_core/theory/reversal_equivalence.tex}.
The single-introduction argument is the idea formalized in
\file{Lean/ConciseIPRA/NativeOneIntro.lean}.
The new construction here is their application to explicitly canonical successor-seeded syntax, with a proof of the resulting change of term equality and substitution.

No independent successor constructor is used. Raw terms are
\[
t::=x\mid 0\mid R(n,a;u.h).
\]
Only \(u\) in \(h\) is bound. Put
\[
Sa:=R(0,a;u.0),\qquad 1:=S0.
\]
Bound variables are renamed to avoid capture. There are no numerical variables over function descriptions: a body is finite earlier syntax, with its remaining free variables as passive parameters.

\section{Canonical terms and substitution}

Define a normalization of raw terms recursively:
\[
\norm{x}=x,\qquad \norm{0}=0,
\]
\[
\norm{R(n,a;u.h)}=
\begin{cases}
R(0,\norm a;u.0),&n\lit0,\\
R(\norm n,\norm a;u.\norm h),&n\not\lit0.
\end{cases}
\]
Here \(n\lit0\) means literal syntactic identity with zero. It is not a numerical equality test. A canonical term is a fixed point of \(\nu\). Equivalently, every subterm with literal zero counter has literal zero body. This is a decidable grammar of finite expressions.

The normalizer preserves the outer term constructor:
\[
\norm t\lit0\quad\Longleftrightarrow\quad t\lit0.
\]
In particular \(\norm{Sa}=S\norm a\). It is idempotent, respects renaming of bound and free variables, and introduces no new free variables.

Extend \(\nu\) to native formulas and contexts by normalizing all terms:
\[
A::=(s=t)\mid Q_y A\mid D_{x,y}A.
\]
The roots \(=\), \(Q\), and \(D\) are preserved. All terms in a canonical formula are canonical.

\begin{lemma}[Substitution coherence]\label{cr:coherence}
For every raw term or formula \(T\) and capture-free simultaneous numerical substitution \(\sigma\),
\[
\norm{T[\sigma]}=
\norm{\,\norm T[\nu\circ\sigma]\,}.
\]
\end{lemma}
\begin{proof}
Use structural induction after freshening bound names. The only additional case is an \(R\)-node. If its original counter is literally zero, both sides erase the body. If it is not literally zero but becomes zero under substitution, both sides again erase the body. In the remaining case both sides retain the \(R\)-node and the equality follows from the induction hypotheses for counter, initializer, and body. The fact that normalization never turns a nonzero-root constructor into literal zero justifies the case division. Formula binders introduce no extra case.
\end{proof}

For canonical expressions define
\[
\subc T\sigma:=\norm{T[\sigma]}.
\]
This is the substitution operation used by every rule below. Lemma~\ref{cr:coherence} gives associativity of normalized substitution and its compatibility with renaming. A variable substituted by zero may expose a zero-counter node. Normalizing only the original syntax, or testing the unsubstituted counter, is insufficient.

In a canonical-calculus display, an expression \(R(n,a;u.h)\) denotes the normalization of that raw constructor expression. This smart-constructor convention is relevant when \(n\lit0\): its written body is discarded. In the raw comparison calculus below, constructor expressions instead retain their literal bodies. Substitution may accept a raw image because it normalizes the completed result; equivalently it may first normalize that image by Lemma~\ref{cr:coherence}.

The sole primitive computation schema is
\[
R(Sn,a;u.h)=\subc h{\norm{R(n,a;u.h)}/u}.
\tag{\(R_+\)}\label{cr:rplus}
\]
All displayed terms are canonical: the left counter \(Sn\) is never literally zero. The raw zero computation
\(R(0,a;u.h)=Sa\) is now definitional, because both sides have the same canonical spelling.

\section{The complete six-rule calculus}

Contexts are finite reusable lists of canonical native formulas. Define root reversal by
\[
(s=t)^\dagger=(t=s),\qquad
(Q_yB)^\dagger=Q_yB,\qquad
(D_{x,y}B)^\dagger=D_{x,y}B.
\]
It commutes with canonical substitution.

The calculus \(\Knorm\) has the following six inference patterns.

\begin{enumerate}[leftmargin=*]
\item \textbf{Hypothesis:} \(\Gamma\vdash A\) when \(A\in\Gamma\).
\item \textbf{Successor computation:} every instance of \eqref{cr:rplus}, in every context.
\item \textbf{Crossed equational transport:}
\[
\frac{\Gamma\vdash s=t\qquad
      \Gamma\vdash \subc E{s/v}}
     {\Gamma\vdash (\subc E{t/v})^\dagger},
\qquad s\not\lit0.
\tag{X}
\]
Here \(E(v)\) is a literal-equation template. Constant templates are allowed.
\item \textbf{Crossed closure elimination:}
\[
\frac{\Gamma\vdash Q_yB(y)\qquad\Gamma\vdash\subc B{0/y}}
     {\Gamma\vdash(\subc B{t/y})^\dagger}.
\tag{QX}
\]
\item \textbf{Paired elimination:}
\[
\frac{\Gamma\vdash D_{x,y}B(x,y)\qquad
      \Gamma,Q_y\subc B{b/x}\vdash A}
     {\Gamma\vdash A}.
\tag{DE}
\]
The eigenparameter \(b\) is fresh for the retained context, conclusion, and fixed parameters of \(B\).
\item \textbf{One introduction:}
\[
\frac{\Gamma,\subc B{t/x,\eta/y}
          \vdash\subc B{t/x,Sz/y}}
     {\Gamma\vdash D_{x,y}B},
\qquad
\eta=
\begin{cases}
z,&B\text{ is a literal equation and }t\lit0,\\
0,&\text{otherwise}.
\end{cases}
\tag{WI}
\]
Here \(z\) is fresh for the retained context, witness \(t\), and fixed parameters of \(B\).
\end{enumerate}

There is no primitive reflexivity, symmetry, numerical case rule, induction rule, zero-computation rule, ordinary logical connective, or general numerical-substitution rule. The designated contradiction is \(0=S0\).

The intended ordinary readings are
\[
(Q_yB)^\circ=B^\circ(0)\to\forall y\,B^\circ(y),
\]
\[
(D_{x,y}B)^\circ=
\exists x\bigl(B^\circ(x,0)\to\forall y\,B^\circ(x,y)\bigr).
\]
All numerical terms use the successor-seeded primitive-recursive interpretation of \(R\). Domains are whole, and equality is actual equality.

\section{First establish the unrestricted comparison}

Do not initially attempt a substitution theorem for the six-rule calculus. Literal-zero source and witness tests are not stable under arbitrary numerical substitution. Instead use two unrestricted comparison calculi.

\begin{definition}[Raw and canonical unrestricted calculi]
The raw calculus \(\Rfree\) uses raw \(R\)-terms, ordinary capture-free substitution, the two computation schemas
\[
R(0,a;u.h)=Sa,\qquad
R(Sn,a;u.h)=h[R(n,a;u.h)/u],
\]
Hyp, unrestricted transport in every native formula template, ordinary QE, DE, and the two introductions
\[
\frac{\Gamma,B(t,0)\vdash B(t,Sz)}
     {\Gamma\vdash D_{x,y}B},
\qquad
\frac{\Gamma,E(t,z)\vdash E(t,Sz)}
     {\Gamma\vdash D_{x,y}E}
\quad(E\text{ a literal equation}).
\]
The canonical calculus \(\Cnorm\) has the same unrestricted rules using canonical substitution, but has only \(R_+\) as its computation schema.

Here ordinary QE is QX without root reversal. Unrestricted transport is X without either its literal-zero source restriction or its root reversal, and it permits arbitrary native formula templates. The freshness conditions for DE and both introductions are the conditions already stated for DE and WI. Hypothesis and DE are unchanged.
\end{definition}

\begin{lemma}[Structural rules]\label{cr:struct}
Weakening, eigenvariable renaming, ordinary-assumption cut, and capture-free numerical substitution are admissible in \(\Rfree\) and \(\Cnorm\).
Weakening, eigenvariable renaming, and ordinary-assumption cut are admissible in \(\Knorm\), without presupposing numerical substitution there.
\end{lemma}
\begin{proof}
Induct on derivations, freshening local eigenparameters first. For \(\Cnorm\), coherence of normalized substitution proves closure of the computation schema and all displayed instances. The case/progress distinction in the unrestricted calculi depends only on the formula root and is stable. Cut replaces assumption uses with proofs without changing witnesses, roots, or source terms; hence it remains valid in the restricted six-rule calculus. Renaming preserves literal zero and root constructors. This proof makes no assertion that general numerical substitution directly preserves an X or WI node.
\end{proof}

\begin{lemma}[Elementary resources]\label{cr:resources}
Both unrestricted calculi derive reflexivity, symmetry, transitivity, numerical congruence, ordinary \(Q\)-introduction by cases, numerical cases, and contextual induction on literal equations.
\end{lemma}
\begin{proof}
Define the retraction context
\[
J(a):=R(S0,0;u.a).
\]
The successor computation gives \(J(a)=a\), even in \(\Cnorm\). Use this equation twice in transport with motive \(v=a\) to get \(a=a\). The usual equality motives give symmetry, transitivity, and numerical congruence.

From \(\Gamma,B(0)\vdash B(Sz)\), case introduction with an unused witness gives \(D_{w,y}B(y)\); DE followed by Hyp exports \(Q_yB(y)\). Applying QE with supplied base and successor proofs gives numerical cases. Repeating the unused-witness construction with progress introduction gives equation induction from its ordinary base and step premises. No selector, predecessor, or ordinary logical connective is needed.
\end{proof}

\begin{lemma}[Recursor body extensionality]\label{cr:body}
Suppose, for a fresh numerical variable \(v\), that a fixed context proves \(h(v)=k(v)\). Then it proves
\[
R(n,a;u.h(u))=R(n,a;u.k(u)).
\]
The assertion holds in both unrestricted calculi.
\end{lemma}
\begin{proof}
Apply contextual literal-equation induction to the equation comparing the two trajectories at a fresh counter. At zero both terms are \(Sa\), by zero computation in \(\Rfree\), and by canonical syntax and reflexivity in \(\Cnorm\). At the successor, computation exposes the two updates. Use the induction equation to replace the old state in \(h\), then use the supplied pointwise body equality at the second old state. Such instantiation is justified by Lemma~\ref{cr:struct}, which has already been proved for these unrestricted calculi. All passive parameters and the initializer remain fixed during the induction.
\end{proof}

This lemma is necessary. Changing a normalized body changes its finite recursor description; it is not generally ordinary congruence of one fixed first-order function symbol.

\begin{lemma}[Normalization is provable in the raw calculus]\label{cr:normal-eq}
For every raw term \(t\),
\[
\Rfree\vdash t=\norm t.
\]
For every raw native formula \(A\), the raw calculus derives both
\(A\vdash\norm A\) and \(\norm A\vdash A\), uniformly in its free parameters.
\end{lemma}
\begin{proof}
For terms use structural induction. Counter and initializer changes use numerical congruence. The body change uses Lemma~\ref{cr:body} and the induction hypothesis with the body's bound variable fresh and free in the auxiliary proof. If the normalized counter is zero, use the raw zero computation to replace the body by zero; otherwise the resulting term is already its normal form.

For formulas use external induction on the native formula syntax. The equation case follows from the two endpoint identities. For the \(Q\) case, under \(Q_yA(y)\), introduce \(Q_y\norm{A(y)}\) by cases: convert its assumed normalized base to the old base, use QE at the fresh successor, and convert back. The reverse direction is identical. For \(D\), eliminate the old \(D\) at a fresh witness, introduce the new \(D\) at that witness by cases, convert the new base to the old base, use the exported old \(Q\), and convert the successor conclusion. Cut closes the construction. These are explicit derived binder-congruence arguments, not applications of an assumed ordinary quantifier logic.
\end{proof}

\begin{theorem}[Canonical and raw consequence]\label{cr:canonical-equiv}
For canonical \(\Gamma,A\),
\[
\Gamma\vdash_{\Cnorm}A
\quad\Longleftrightarrow\quad
\Gamma\vdash_{\Rfree}A.
\]
More generally, normalization sends every raw derivation to a canonical derivation of its normalized sequent.
\end{theorem}
\begin{proof}
For the raw-to-canonical direction normalize the entire finite derivation. Coherence of substitution handles every instance of transport, QE, DE, and either introduction. A raw zero-computation leaf normalizes to \(Sa=Sa\), which is derived by Lemma~\ref{cr:resources}; a successor-computation leaf normalizes to an \(R_+\) instance.

For the converse, regard canonical expressions as raw expressions and recurse on the canonical derivation. Any canonical rule instance differs from its raw counterpart only by normalizations of substituted expressions. Lemma~\ref{cr:normal-eq} converts these premises, local assumptions, and conclusions in both directions. Use ordinary-assumption cut for local assumptions. For example, canonical transport first converts the canonical minor \(\nu(E[s/v])\) to raw \(E[s/v]\), applies raw transport, and converts \(E[t/v]\) to its canonical normalization. Introductions and DE use the same conversions under their retained local assumptions. Freshness is preserved because normalization introduces no free variables. No restricted-rule substitution theorem enters either direction.
\end{proof}

\section{Only then reduce the canonical calculus to six rules}

\subsection{Equality before introduction repair}

In \(\Knorm\), keep \(J(a)=R(S0,0;u.a)\) and \(c=J(0)\). Each \(J(a)\) is literally nonzero and \(J(a)=a\) is an \(R_+\) instance.

X with major \(c=0\) and a constant equation template reverses any proved equation. Compound root reversal is the identity. Therefore QX followed by reversal supplies ordinary QE, and X followed by reversal supplies ordinary equational transport at nonzero literal sources.

Reflexivity follows by using \(J(a)=a\) as both X premises with template \(v=a\). To remove the literal-zero source restriction, use the established retraction argument, whose short details are included here.

If \(p,q\not\lit0\), \(p=0\), and \(q=0\), then X with major \(p=0\), minor \(J(p)=p\), and template \(J(v)=p\) yields \(p=c\). Similarly \(q=c\). Reverse and join through the allowed source \(c\) to obtain \(p=q\). Transitivity now holds without restriction: if its middle term is nonzero use restricted ordinary transport; if it is zero and both endpoints are nonzero, reverse the second equation and use the zero join; literal-zero endpoint cases are immediate.

Numerical congruence along a nonzero source is ordinary restricted transport from reflexivity. For an equation \(0=t\) with \(t\not\lit0\), reverse it to \(t=0\); transport the reflexivity \(f(t)=f(t)\) in the template \(f(v)=f(t)\), giving \(f(0)=f(t)\). If both endpoints are zero, use reflexivity. Finally, from \(s=t\) and \(a(s)=b(s)\), numerical congruence and unrestricted transitivity give
\[
a(t)=a(s)=b(s)=b(t).
\]
Thus unrestricted ordinary transport in literal-equation templates is available before any introduction repair. All steps instantiate explicit schemas at the terms at hand; none substitutes into a restricted derivation.

\subsection{Recovering both introductions}

\begin{lemma}[Case introduction]\label{cr:case}
Every unrestricted case-introduction instance is derivable in \(\Knorm\).
\end{lemma}
\begin{proof}
WI already is case introduction when its matrix is compound or its witness is not literally zero. Only an equation matrix \(E\) at witness zero remains.

Use instead the witness \(c=J(0)\), which is nonliteral and provably zero. Under the local assumption \(E(c,0)\), unrestricted equation transport gives \(E(0,0)\). Cut in the given case-premise derivation to get \(E(0,Sz)\); transport back to \(E(c,Sz)\). WI at witness \(c\) is the case choice and concludes the required unchanged formula \(D_{x,y}E(x,y)\). The new witness contains no eigenparameter \(z\).
\end{proof}

\begin{lemma}[Progress introduction]\label{cr:progress}
Every unrestricted literal-equation progress instance is derivable in \(\Knorm\).
\end{lemma}
\begin{proof}
Suppose the supplied premise is
\(\Gamma,E(t,z)\vdash E(t,Sz)\).
Let the first coordinate \(w\) be unused in the equation \(E(t,y)\). WI with witness zero proves
\[
\Gamma\vdash D_{w,y}E(t,y)
\]
by its progress branch. Under the additional assumption \(E(t,0)\), eliminate this paired formula. Its exported clause is \(Q_yE(t,y)\), independent of the fresh witness. Ordinary QE gives \(E(t,Sz)\). Closing DE gives the premise required for the recovered case introduction at the original witness \(t\), which concludes \(D_{x,y}E(x,y)\).
\end{proof}

\subsection{Compound transport}

Now equation transport, ordinary QE, DE, and both introductions are available. Derive ordinary transport in all native formula templates by external induction on their constructor structure. The equality case is already proved. For a \(Q\)-template, use derived case \(Q\)-introduction: convert its new base back along the equality by the induction hypothesis, apply the old \(Q\) at the fresh successor, and convert forward. For a \(D\)-template, DE chooses the old witness; case-introduce the new \(D\) at that witness, convert its base back, use the exported old clause, and convert the successor forward. Reverse the numerical equality to obtain both directions when required. The derivations use only contextual equation symmetry, cut, and renaming already established above.

\begin{theorem}[Six-rule equivalence]\label{cr:six}
For every canonical native sequent,
\[
\Gamma\vdash_{\Knorm}A
\quad\Longleftrightarrow\quad
\Gamma\vdash_{\Cnorm}A.
\]
\end{theorem}
\begin{proof}
The preceding constructions provide every rule of \(\Cnorm\) inside \(\Knorm\). Conversely, a WI instance is exactly one of the two unrestricted introductions, and the crossed rules are ordinary transport or QE followed by derived equality symmetry when their conclusion is an equation. Hyp, \(R_+\), and DE are common. Transform finite derivations recursively. No numerical-substitution admissibility for \(\Knorm\) was presupposed; it may now be transferred through the equivalence.
\end{proof}

\paragraph{Syntactic tests under normalization.}
Normalization preserves literal-zero status of a whole term and the root constructor of a formula, so it does not itself alter an X or WI test. Arbitrary numerical substitution can alter such tests. Keeping the unrestricted equivalence and the later restricted-rule repair separate avoids confusing these two facts.

\section{Arithmetic bridge and both interpretation composites}

The full arithmetic bootstrap for \(\Rfree\) is inherited, not a new claim of this note. Appendices~\ref{nb:section} and~\ref{js:section} reproduce the pinned native finite-case repair and successor-seeded predecessor construction. Their recursor is denoted \(J\) rather than \(R\); renaming that one constructor makes it the present raw language.

Their explicit term \(P\) is proved, with only native cases, equality reasoning, and literal-equation induction, to satisfy
\[
P0=0,\qquad P(Sa)=a.
\]
No predecessor, successor cancellation, numerical equality test, ordinary selector, or compound-formula induction is assumed during that bootstrap. The native finite-case appendix is required: theorem-level numerical cases alone would not justify arbitrary contextual finite case splits before the logical macros exist.

The initialized scalar iterator is then the finite raw \(R\)-term
\[
\widehat I(n,a;u.h)
=
P\!\left(R(n,a;v.S\bigl(h[P(v)/u]\bigr))\right).
\]
Its zero computation is \(P(Sa)=a\). At a successor, \(R_+\) exposes an outer successor that \(P\) removes, leaving exactly
\[
\widehat I(Sn,a;u.h)=h[\widehat I(n,a;u.h)/u].
\]
The unrestricted selector \(C(n,a,b):=\widehat I(n,a;u.b)\) is now available. The established native scalar bridge reconstructs all intuitionistic first-order connectives, actual-equality reflection, quantifier-free induction, pairing, and every fixed primitive-recursive definition. That bridge is separately formalized in the pinned repository; its validity is inherited here.

Conversely, interpret a raw successor-seeded description by
\[
R(n,a;u.h)\longmapsto I(n,Sa;u.h).
\]
The two raw computations are instances of the scalar computations. For the two numerical interpretation composites, first identify successor: its definition \(R(0,a;u.0)\) translates to \(I(0,Sa;u.0)=Sa\). Then proceed through each fixed description in dependency order. After earlier descriptions have been identified, the description and its double translation have the same initial value and the same update. Literal-equation induction on the equation comparing their trajectories proves equality throughout the entire numerical domain. The initializer and passive parameters remain independent and fixed during this induction, and are instantiated only afterward. Body equality is used through the derived trajectory comparison, not as an uncharged recursor-congruence axiom.

The native formula round trips use the established logical macros once \(C\) is available. The \(Q\) and \(D\) readings are those already displayed; their inverse translations are interderivable by case introduction, DE, and QE. Formula congruence and the numerical identities extend the round trips to every formula.

Combining these inherited arithmetic and logical results with Theorems~\ref{cr:canonical-equiv} and~\ref{cr:six} gives the precise conclusion:
\[
\Knorm\ \text{is equiconsistent and bi-interpretable with full first-order }\ipra.
\]
Both numerical domains are whole and interpreted equality is actual equality. The numerical composite isomorphisms have graph \(x=y\). On raw syntax the composite through canonical syntax is provably, rather than literally, the identity; on canonical syntax normalization is literally idempotent.

\section{What has and has not become simpler}

There are six explicit inference patterns instead of seven, and no independent successor constructor. The omitted zero computation has become a canonical-syntax law:
\[
R(0,a;u.h)\equiv R(0,a;u.0).
\]
That law must be enforced by every substitution and comparison used by a complete checker. It has not disappeared from the trusted specification.

For the current compact encodings, a numerical recursor and a \(Q\)-formula both have arity three, so their opener must distinguish numerical depth before applying this normalization. Testing the original counter instead of the substituted one is unsound. Merely rejecting noncanonical input also fails: internally generated substitution instances must be normalized.

Accordingly this theorem establishes a mathematical change of presentation and a smaller explicit rule inventory. The separately measured matching strict implementation is also smaller; no absolute or compact-checker minimum for this new calculus follows. A fair measurement must include canonicalization, syntax and scope handling, decoding, every inference rule, and the public input/output driver. Any separate implementation improvement for the old scalar calculus must not be attributed to this new canonical \(R\) calculus.



\section{Complete implementations and measured complexity}
\label{sec:implementation}

\subsection{Matching strict implementations}

The directory \file{strict/} contains three standalone Python modules.
Each imports only \texttt{json} and \texttt{sys}; none imports a proof
producer, axiom database, normalization oracle, or test utility. The public
input is a JSON object with exactly \texttt{goal} and \texttt{proof} fields.
The goal must be closed and satisfy the entire stated grammar. A successful
run exits with status zero and prints exactly \texttt{ACCEPT}; every failed
check rejects.

The three modules use the same proof-tree format, the same context handling,
the same de Bruijn convention, the same logical rules, and the same public
driver. The scalar control has explicit successor and initialized iteration.
It uses crossed eliminations and witness-directed introduction, so it is an
equivalent matching presentation of the scalar theory, rather than a literal
transcription of the uncrossed \texttt{MPrf} declaration.
The raw-$R$ control replaces successor by its recursor abbreviation and keeps
both computations. The canonical-$R$ module additionally normalizes every
recursor construction, validates canonical inputs, and deletes the explicit
zero-computation rule. This is a comparison of complete matching
implementations, not a comparison of a fragment against a complete module.

\begin{center}
\small
\begin{tabular}{lrrrr}
\toprule
Complete module & Tokens & AST & Bytes & Gzip 9\\
\midrule
Scalar $I$, seven rules & 1433 & 1595 & 4923 & 1224\\
Raw $R$, seven rules & 1378 & 1532 & 4753 & 1209\\
Canonical $R$, six rules & 1338 & 1485 & 4611 & 1205\\
\bottomrule
\end{tabular}
\par\smallskip
Matching strict implementations; normalization and both interfaces are charged.
\end{center}


The canonical program saves 312 normalized bytes against the scalar control,
or $6.34\%$, together with 95 significant tokens and 110 AST nodes. Against
the raw-$R$ control it saves 142 normalized bytes. Thus normalization has
been charged, and the resulting strict implementation is still smaller in
this matched design. These programs are intentionally transparent; their
size does not improve the repository's much more compressed arity codec.

Measurements use CPython 3.12.14. Actual docstrings are removed, the complete
module is normalized with \texttt{ast.unparse}, and one final LF is appended.
Significant tokens exclude comments, indentation, dedentation, and line-end
and end-marker tokens. AST nodes are counted with \texttt{ast.walk}.
The last column is gzip level 9 of the normalized module with timestamp zero;
it is an informational coordinate, not the optimization objective. Raw
source hashes and sizes are recorded separately.

\subsection{Checking syntax, normalization, and binders}

In the strict format, $0$ is JSON integer zero, positive integers are
de Bruijn variable indices, and $[n,a,h]$ is an $R$ term whose body binds
index one. Formula tags are \texttt{=}, \texttt{Q}, and \texttt{D}. In a
$D$ matrix, the universal variable has index one and the existential
variable index two. Booleans, floats, negative indices, unknown tags, and
extra fields are rejected where the specification requires an integer,
formula, or fixed rule record.

The smart constructor is the complete normalization operation at a newly
constructed recursor node:
\[
 \operatorname{rec}(n,a,h)=[n,a,\;0\text{ if }n=0\text{ else }h].
\]
Its children have already been normalized by recursive substitution.
Counter and initializer are traversed at the current binding depth; the body
is traversed one binder deeper. Substituted images are shifted under binders.
The implementation therefore normalizes a zero counter newly exposed by
substitution, and it never mistakes a bound state for a free parameter.

The strict test suite contains 86 named certificates across the three
calculi. All 172 command-line replays of the readable and normalized modules
match their expectations. The canonical calculus has 12 positive and 18
negative certificates. Its positive cases include each inference mechanism,
an actually used induction hypothesis, a paired elimination with a genuine
witness and shifted retained context, and zero-counter normalization exposed
by transport. Negative controls include a future witness, an escaping
eigenparameter, an invalid source, a wrong closure base, a compound transport
template, malformed syntax, and use of the removed zero rule.

There are also 480 independent numerical evaluations comparing substitution
with environment evaluation, plus direct tests of nesting and capture.
These are finite executable controls. The universal soundness and
completeness argument is the correspondence between the six checked rule
cases and the calculus proved equivalent above; the numerical tests alone
do not establish that metatheorem.

\subsection{Separate compact implementation refinement}
\label{sec:compact-refinement}

The archive also contains a small refinement of the \emph{existing} guarded
scalar calculus. It has no role in the canonical-$R$ equivalence theorem.
The source transformations pair the already checked major and target
comparisons, use a shorter integer/list discriminator, and compose the
existing scalar-zero administrative truth convention with the guarded
checker. A conservative variant changes only the boolean spelling of the
expected major and preserves all JSON executions of its parent.

\begin{center}
\small
\begin{tabular}{lrrrr}
\toprule
Complete module & Tokens & AST & Bytes & Gzip 9\\
\midrule
Frozen compact byte control & 265 & 303 & 522 & 329\\
Refined compact byte profile & 265 & 306 & 516 & 332\\
Frozen compact AST control & 262 & 281 & 606 & 348\\
Refined compact AST profile & 263 & 280 & 602 & 348\\
Matched strict guarded control & 541 & 620 & 1381 & 633\\
Refined strict guarded checker & 542 & 624 & 1376 & 633\\
\bottomrule
\end{tabular}
\par\smallskip
Existing scalar calculus only. Byte and AST improvements use different modules.
\end{center}


The compact byte and AST coordinates belong to different programs. Their
public-language qualification remains the inherited guarded Eq/$D$ contract:
the compact programs do not themselves validate the public grammar. The
directory includes a separately counted strict goal and integer/list
certificate front end, frozen controls, converted fixtures, an exact
description of the changed neutral marker, and a reproducible measurement
and test script. The compact gain is an implementation result, not a new
arithmetic basis and not the size of the canonical-$R$ checker.

The eager paired comparison may inspect previously unused instruction
fields. Its completeness statement is therefore made for the specified
well-scoped integer/list certificate language, and no equality of arbitrary
malformed-JSON behavior is claimed for that variant. Full statements and
the execution evidence are in the compact directory. No failed experimental
checker is supplied as an accepted implementation.



\section{Why deleting the initializer is a different problem}
\label{sec:scalar-obstruction}

One might try to remove the initializer argument altogether. This section
gives a new scalar-specific obstruction to that route. The argument is about
standard term functions, not about a lower bound for all possible
axiomatizations or proof checkers.

Consider terms generated by variables, $0$, $S$, and the scalar constructor
\[
 F(0;v.h)=S(h(0)),\qquad
 F(Sn;v.h)=h(F(n;v.h)).
\]
Arbitrary nesting and passive numerical parameters are allowed. There is no
primitive numerical equality characteristic or additional arithmetic
operation. The successor may instead be the abbreviation $F(0;v.a)$.

\begin{theorem}[A restriction on every unary scalar term]
\label{thm:unary-bound}
For every unary term $t(n)$ in this language, there is a natural number $C$
such that
\[
 \forall n\in\mathbb N\quad
       \bigl(t(n)\leq C\ \lor\ t(n)\geq n\bigr).
\]
Consequently no such term defines predecessor $\max(n-1,0)$, integer
halving, or any function having unbounded values strictly below its input.
\end{theorem}

The counter-aware tied-seed language in the existing
\texttt{local\_scopes} package already has a finite-control obstruction, but
it does define predecessor. Thus the unary conclusion above is specific to
removing access to the old counter as well as the independent initializer.
We strengthen the existing finite-control invariant; its original complete
proof is included in the archive as
\file{reference/finite_profiles.tex}.

\begin{definition}
Let $\mathcal U$ be the unary functions $f:\mathbb N\to\mathbb N$ for which
there is $C$ satisfying $f(n)\leq C$ or $f(n)\geq n$ at every input.
\end{definition}

\begin{lemma}[Closure properties]
Constants, the identity, and successor belong to $\mathcal U$. The class is
closed under composition and under pointwise selection from a fixed finite
family of its members.
\end{lemma}
\begin{proof}
For finite selection, take the maximum of the finitely many bounds. For
composition $f\circ g$, fix bounds $C_f,C_g$. If $g(n)\leq C_g$, then the
output is bounded by $\max_{j\leq C_g}f(j)$. Otherwise $g(n)\geq n$, and
the output is either at most $C_f$ or at least $g(n)\geq n$.
\end{proof}

\begin{lemma}[The orbit lemma]
If $h\in\mathcal U$, then the sequence
\[
 r(0)=1+h(0),\qquad r(n+1)=h(r(n))
\]
viewed as a unary function of $n$, also belongs to $\mathcal U$.
\end{lemma}
\begin{proof}
A finite orbit is bounded. An infinite orbit has no repeated value, since
a repeat in a deterministic orbit makes its future periodic. Fix a bound
$C$ for $h$. A strictly decreasing transition must end in
$\{0,\ldots,C\}$, so an injective orbit has only finitely many such
transitions. After the last one, the orbit increases strictly. Eventually
it exceeds every value in the earlier finite prefix. For all sufficiently
large $n$, its first $n+1$ distinct values therefore lie in
$[0,r(n)]$, and $r(n)\geq n$. Increase a constant to cover the finitely many
earlier outputs.
\end{proof}

\begin{proof}[Proof of Theorem~\ref{thm:unary-bound}]
Strengthen the finite-control invariant as follows. Every positive-arity
term has one active numerical coordinate; each other coordinate is observed
through a finite unary feature map; once those finite observations are fixed,
the output is a unary profile of the active coordinate. Require every such
profile to belong to $\mathcal U$. Constants and projections satisfy this
property. In a composition, fixing the inactive features leaves a pointwise
finite selection of compositions of old profiles. The closure lemma applies.

For $F$, first treat its counter as an independent input. If the active
coordinate of its body is the state $v$, fixing the finitely many passive
features gives finitely many updates $h_c\in\mathcal U$. Their orbit profiles
belong to $\mathcal U$ by the orbit lemma, and the counter is the new active
coordinate.

If instead the active coordinate of the body is a passive parameter $p$,
write
\[
 h(v,p,\vec q)=H_{\beta(v),\gamma(\vec q)}(p),
\]
where $\beta,\gamma$ take finitely many values and every $H_{b,c}$ belongs
to $\mathcal U$. Form a finite catalogue $V_j(p)$ of all update profiles
$H_{b,c}(p)$ and all initial profiles $1+H_{\beta(0),c}(p)$. Every orbit
state is a catalogue value, and every catalogue member belongs to
$\mathcal U$.

The finite feature
\[
 \kappa(p)=(\beta(V_j(p)))_j
\]
determines every transition of the catalogue indices once
$c=\gamma(\vec q)$ is fixed. For every possible $c,k$, let
$\theta_{c,k}(n)$ be that finite-state process after $n$ updates. Pack all
these finitely many processes into the finite unary feature
$\eta(n)=(\theta_{c,k}(n))_{c,k}$. Then
\[
 F(n;v.h)=V_{\eta(n)_{\gamma(\vec q),\kappa(p)}}(p).
\]
For fixed features of $n,\vec q$, this is a pointwise finite selection of
the catalogue profiles of $p$, so it belongs to $\mathcal U$. This argument
uses only the displayed finite observations, not a numerical equality
oracle. Substitute the actual counter term by the already established
composition closure. Structural induction is complete. At arity one there
are no inactive coordinates, yielding the asserted unary conclusion.
\end{proof}

For predecessor, choose $n>C+1$; its value is then greater than $C$ and less
than $n$, a contradiction. Halving fails at arbitrarily large $n>2C+1$.
The theorem by itself concerns term definitions. It must not be presented
as an obstruction to arbitrary graph interpretations or changed domains.
Those broader questions are treated separately in the existing relational
obstruction, also supplied in the archive.

\subsection{Independence of the original two scalar computations}

Neither computation equation of the original initialized scalar iterator
can simply be deleted, even if full induction is supplied. In the standard
natural numbers, interpret every fixed body compositionally.

If $I_0$ is omitted, put $I(n,a;h)=h^n(0)$. The successor equation holds,
but $I(0,1;u.u)=0\ne1$. If $I_S$ is omitted, put $I(n,a;h)=a$. The zero
equation holds, but the successor equation fails for $h(u)=Su$.
Both structures have ordinary numerical coverage, separation, and every
first-order induction instance. Thus the canonical-$R$ construction succeeds
by changing the representation; it does not derive a deleted scalar law
from the unchanged remaining axioms.



% Includable mathematical section; requires amsmath, amssymb, amsthm.
% Requires theorem, lemma, corollary, proposition environments.
% New countermodel, independently reviewed during this request.
% This section proves failure of full equation-induction recovery.
% It does NOT prove absence of every possible bi-interpretation with iPRA.

\section{Zero-equation induction does not recover equation induction}
\label{zi:section}

Consider the initialized scalar language
\[
t::=x\mid0\mid St\mid I(n,a;u.h)
\]
with its two computation schemas
\[
I(0,a;u.h)=a,\qquad
I(Sn,a;u.h)=h[I(n,a;u.h)/u].
\]
Only the state variable \(u\) is bound in \(h\); the other variables of the fixed finite body are numerical parameters. Use ordinary intuitionistic first-order logic with actual equality, numerical coverage
\[
\forall x\,(x=0\lor\exists y\,x=Sy),
\]
separation \(0=S0\to\bot\), and contextual induction restricted to equations
\[
0=t(n,\vec p)
\]
(or, equivalently, \(t(n,\vec p)=0\)). Call this theory \(T_0\).
Adding successor injectivity will not affect the counterexample.

A tempting simplification is to use progressive introduction only when a matrix equation has literal zero on one side, and to use the case introduction otherwise. The following theorem rules out the intended recovery of the existing full equation-induction rule from this restriction.

\begin{theorem}[A model separating zero induction from equation induction]
\label{zi:main}
There is a classical model of \(T_0\), also satisfying successor injectivity, in which
\[
\forall n\,\forall a\; I(n,a;u.u)=a
\]
is false. Nevertheless its zero instance and its ordinary induction step are theorems of the computation equations and equality logic.
Consequently induction on arbitrary literal equations is not admissible in \(T_0\).
\end{theorem}

The model is a classical one-world model and hence also a model for intuitionistic first-order logic. Its construction takes place in ordinary classical mathematics, using compactness and choices of compatible paths. These are metatheoretic resources for constructing a countermodel, not extra principles granted to \(T_0\).

\subsection{An arithmetic base with two-point fibers}

Fix a nonstandard model \(M\) of full Peano arithmetic expanded by the usual symbols for all fixed primitive-recursive descriptions. Such a model exists in ordinary mathematics: add a constant larger than each standard numeral, use the standard natural-number model for every finite subtheory, and apply first-order compactness. A countable model may be chosen. We use its full induction only in the external countermodel construction.

For \(m\in M\), define the finite nonempty set
\[
A_m=
\begin{cases}
\{0\},&m\text{ is a standard numeral},\\
\{0,1\},&m\text{ is nonstandard}.
\end{cases}
\]
The distinction is external; it need not be definable in \(M\).
Let
\[
D:=\{(m,i):m\in M,\ i\in A_m\},\qquad
\pi(m,i):=m.
\]
Define
\[
0_D=(0_M,0),\qquad
S_D(m,i)=(S_Mm,i).
\]
This is well defined because successor preserves standardness and nonstandardness.

The map \(\pi\) is surjective and preserves zero and successor. Every standard numeral has a singleton fiber; every nonstandard element has two lifts. Successor is injective on \(D\), and zero is outside its image. Coverage holds: a positive standard numeral has its usual predecessor, and a nonstandard \((m,i)\) has predecessor \((m-1,i)\).

The nonstandard part of \(M\), under successor and predecessor, is a disjoint union of externally bi-infinite chains. More precisely, each equivalence class under a difference of a standard integer is a copy of \(\mathbb Z\). On \(D\), the label \(i\) is retained along each such counter chain.

\subsection{Compatible paths through finite fibers}

\begin{lemma}[Finite-fiber path lemma]\label{zi:path}
Let \((B_j)_{j\in\mathbb Z}\) be finite nonempty sets and let
\(f_j:B_j\to B_{j+1}\) be total maps. There is a sequence
\((b_j)_{j\in\mathbb Z}\) such that
\[
b_j\in B_j,\qquad b_{j+1}=f_j(b_j)
\]
for every \(j\).
\end{lemma}
\begin{proof}
The product \(\prod_{j\in\mathbb Z}B_j\), with the finite discrete factors, is compact. Each adjacent-coordinate equation is a closed constraint. Every finite collection of these constraints is satisfiable: choose a value at the least involved index and propagate it forward, choosing arbitrary values elsewhere. The finite-intersection property therefore gives a simultaneous solution. Equivalently, this is the usual compactness argument for paths in a sequence of finite nonempty layers. Surjectivity of the transition maps is not required.
\end{proof}

\subsection{Choosing iterators by actual body functions}

It is essential not to make independent choices merely for syntactic descriptions of bodies. Two bodies that denote the same actual function must receive the same iterator. The following construction ensures this and makes ordinary substitution under the binder legitimate.

Let \(\mathscr H\) be the set of all actual functions \(h:D\to D\) for which there is a unary primitive-recursive function with parameters \(f:M\to M\) satisfying
\[
\pi(h(x))=f(\pi(x))\qquad(x\in D).
\tag{Projection}\label{zi:projection}
\]
Surjectivity of \(\pi\) makes \(f\) unique as an actual function on \(M\).

For such an \(f\), write \(\operatorname{It}_f(m,b)\) for its usual iteration in \(M\), with initial value \(b\). This notation does not depend on the chosen primitive-recursive description and parameter tuple for \(f\). Indeed, if two descriptions give the same function on all of \(M\), their two iterators have the same base value and pointwise identical updates; the equation-induction instance in \(M\) comparing them proves equality at every counter. This is a use of induction in the external base model, not in the weak theory \(T_0\).

We define one external operation
\[
\mathcal I:D\times D\times\mathscr H\longrightarrow D.
\]
It will obey the scalar computation equations and satisfy
\[
\pi(\mathcal I(n,a,h))
=\operatorname{It}_f(\pi(n),\pi(a)).
\tag{Iterator projection}\label{zi:it-projection}
\]

If \(n\) is a standard counter \(k\), put
\[
\mathcal I(k,a,h):=h^k(a),
\]
where the right side is actual finite iteration on \(D\).

Now fix \(a,h\) and one nonstandard counter chain
\[
n_j=(m_j,i),\qquad
m_{j+1}=S_Mm_j,\qquad j\in\mathbb Z.
\]
Its counter label \(i\) is fixed. Put
\[
r_j:=\operatorname{It}_f(m_j,\pi(a)),\qquad
B_j:=\pi^{-1}(r_j).
\]
The sets \(B_j\) have one or two elements. By the computation law in \(M\),
\(r_{j+1}=f(r_j)\). Equation~\eqref{zi:projection} therefore makes the restriction of \(h\) a total map
\[
h:B_j\longrightarrow B_{j+1}.
\]
Lemma~\ref{zi:path} gives a compatible path \(b_j\in B_j\). Choose one such path and define
\[
\mathcal I(n_j,a,h):=b_j.
\]
Do this for every triple consisting of an actual function \(h\), an initializer \(a\), and a nonstandard counter chain. The choices may be made simultaneously in the metatheory.

There is one specified family of choices. If \(h=\operatorname{id}_D\), choose on every nonstandard counter chain the constant path
\[
\mathcal I(n,a,\operatorname{id}_D):=(\pi(a),0).
\tag{Identity choice}\label{zi:identity-choice}
\]
This is legitimate: the projected iterator is constantly \(\pi(a)\), the label-zero point always exists, and the identity update fixes it. These choices depend on the actual function \(h\), so every syntactic presentation of the identity body uses the same family.

By construction,
\[
\mathcal I(0_D,a,h)=a,\qquad
\mathcal I(S_Dn,a,h)=h(\mathcal I(n,a,h))
\]
for every \(n,a,h\), and \eqref{zi:it-projection} holds. The standard and nonstandard counter components never meet under a finite successor step, so their definitions are compatible.

\subsection{Interpreting every finite scalar term}

Interpret the binding constructor by
\[
\llbracket I(n,a;u.h)\rrbracket_D
=
\mathcal I\bigl(\llbracket n\rrbracket_D,
              \llbracket a\rrbracket_D,
              v\mapsto\llbracket h\rrbracket_{D,u:=v}\bigr).
\]

\begin{lemma}[Term projection and absence of syntactic circularity]
\label{zi:term-projection}
Every finite scalar term has a well-defined interpretation on \(D\). For every valuation \(\rho\),
\[
\pi(\llbracket t\rrbracket_{D,\rho})
=
\llbracket t\rrbracket_{M,\pi\circ\rho},
\]
where the right side uses the ordinary initialized scalar iterator in \(M\). In particular every body function supplied to \(\mathcal I\) belongs to \(\mathscr H\).
\end{lemma}
\begin{proof}
Induct externally on the finite syntax, treating a body's bound variable as an additional free numerical input in the induction hypothesis. Zero, successor, and variables are immediate. The induction hypothesis for a body gives a function satisfying \eqref{zi:projection}, whose projected function is the fixed primitive-recursive interpretation of that body, with its fixed parameter values in \(M\). Thus it lies in \(\mathscr H\), and \eqref{zi:it-projection} proves the assertion for the enclosing iterator.

The operation \(\mathcal I\) was chosen for actual functions in \(\mathscr H\) before this syntax induction; its definition does not call an evaluator for the term language. Finite nesting therefore creates no circular semantic definition. Equivalently, one obtains a first-order structure by taking one function symbol for every finite earlier-stage body and giving its numerical arguments the interpretation displayed above. Equal parameter values determine the same actual body function and hence the same iterator. The usual substitution lemma, including substitution under a recursor body after renaming bound variables, follows from compositional semantics.
\end{proof}

This also shows that the two scalar computation schemas hold for every fixed finite body, including nested descriptions. No uniform total evaluator is postulated in the object language.

\subsection{Why every zero-equation induction instance holds}

A singleton zero fiber gives
\[
\llbracket t\rrbracket_D=0_D
\quad\Longleftrightarrow\quad
\llbracket t\rrbracket_M=0_M
\]
after projecting the valuation.

Fix numerical parameters \(\vec a\in D\) and a scalar term \(t(n,\vec a)\). Suppose in \(D\) that
\[
t(0,\vec a)=0,\qquad
\forall n\,\bigl(t(n,\vec a)=0\to t(Sn,\vec a)=0\bigr).
\]
Write \(\overline t\) for the primitive-recursive interpretation in \(M\).
Term projection and surjectivity of \(\pi\) imply
\[
\overline t(0,\pi\vec a)=0,\qquad
\forall m\in M\,
\bigl(\overline t(m,\pi\vec a)=0
      \to\overline t(Sm,\pi\vec a)=0\bigr).
\]
Apply the relevant induction instance in \(M\), then reflect zero back through \(\pi\), to obtain \(\forall n\in D\,t(n,\vec a)=0\).

This validates contextual zero-equation induction as well. At any valuation where the retained context is true, the two premises supply exactly the displayed base and universal step; no projection of the possibly complicated retained context is required. Thus arbitrary first-order side assumptions do not repair the missing induction strength.

In fact the same argument validates induction on every equation
\[
\overline{k}=t(n,\vec a)
\]
with one side a fixed standard numeral, because every such numeral has a singleton fiber. Allowing every fixed numeral in place of zero still does not recover full equation induction.

\subsection{Failure of identity-iterator uniqueness}

Choose a nonstandard \(b\in M\) and let \(a=(b,1)\). For the identity body define
\[
E(n):\quad I(n,a;u.u)=a.
\]
Its base is a computation equation. Its step is provable using
\[
I(Sn,a;u.u)=I(n,a;u.u):
\]
under \(E(n)\), equality transitivity gives \(E(Sn)\). Hence these are genuinely provable induction premises, not merely premises accidentally true in the constructed model.

For every nonstandard counter \(n\), however, \eqref{zi:identity-choice} gives
\[
I(n,a;u.u)=(b,0)\ne(b,1)=a.
\]
This proves Theorem~\ref{zi:main}. The failed equality is between two elements with the same projection; zero tests cannot distinguish the two labels.

\begin{corollary}[No recovery of actual-equality reflection]
\label{zi:no-reflection}
The theory \(T_0\) cannot prove, for a scalar term \(e(x,y)\), the reflection law
\[
e(x,y)=0\quad\longleftrightarrow\quad x=y.
\]
\end{corollary}
\begin{proof}
Take the distinct points \(x=(b,0)\), \(y=(b,1)\) over one nonstandard \(b\). Term projection gives the same projected value to \(e(x,y)\) and \(e(x,x)\). If the proposed reflection law held, \(e(x,x)=0\); the singleton zero fiber would then force \(e(x,y)=0\) also, contradicting \(x\ne y\).
\end{proof}

The usual conversion of arbitrary equations into zero equations therefore cannot be used to justify this induction restriction: in this scalar presentation the required actual-equality reflection is precisely a missing conclusion.

\paragraph{Exact scope.}
The theorem refutes recovery of the current scalar equation-induction rule, even with full intuitionistic first-order equality logic, all numerical cases, both computation schemas, and successor injectivity available. In the native \(Q,D\) presentation, the ordinary readings make QE, DE, case introduction, and progress for zero equations valid in this same model; hence those logical rules do not evade the counterexample.

The construction does not prove that \(T_0\) has no interpretation of arithmetic on a quotient, and does not prove that every possible changed-domain bi-interpretation with \(\ipra\) is impossible. Such claims require additional arguments. Its precise conclusion is that the proposed literal-zero progress restriction does not preserve the existing same-domain, actual-equality consequence relation.

\paragraph{Further research suggested by the countermodel.}
A useful remaining question is to characterize which additional zero-induction principles force the fibers invisible to zero tests to collapse. Actual-equality reflection suffices but is a substantial additional resource. Another candidate is a narrowly stated iterator-uniqueness or extensionality principle that excludes the identity-iterator defect above; its sufficiency for all equation induction, and its complete-checker cost, remain to be proved.



\section{Further research questions}

\begin{enumerate}[leftmargin=*]
\item Can a canonical-$R$ encoding absorb zero-counter erasure into an
already necessary traversal and improve the current compact source minima,
while validating its complete public language? The normalization must be
applied after substitution, not merely when reading the original input.
\item Can the present normalization equivalence be formalized by extending
the existing Lean term and native-calculus development? Body extensionality
should be proved before normalization soundness; literal-source and
witness tests should be handled after unrestricted equivalence.
\item Which additional principle, beyond zero-equation induction, rules out
the two-point fibers of Section~\ref{zi:section}? Actual-equality reflection
suffices, but its cost must be counted. Does a narrowly stated
identity-iterator uniqueness principle already suffice for full equation
induction?
\item Can the new unary obstruction be strengthened to a useful normal
form or effective classification of the scalar self-seeded term functions?
The proof above supplies existence of a bound, not an algorithm extracting
an optimal bound from a term.
\item Can a pointfree proof representation reduce the substitution
machinery without introducing arbitrary extra pairing structure? External
tuple contexts avoid that interpretation-composite problem, but their
typing and conversion costs must be counted.
\end{enumerate}

\paragraph{Research status.}
The six-rule theorem, unary obstruction, and zero-induction countermodel are written mathematical
arguments. The executable checks test the stated implementations, scoping,
normalization, and selected derivations; they do not constitute a formal
verification of either universal metatheorem. The upstream formalized
scalar bridge is used with its stated scope. No global minimality or
literature-wide priority theorem is asserted.



\clearpage\appendix

\section*{Attribution for the inherited appendices}
The next seven source sections are reproduced from the pinned ProveIt
snapshot under its MIT No Attribution license. They are reference material,
not new contributions of this report. Their original notations and internal
attribution are retained. The two successor-seeded sections justify the raw
$R$ arithmetic used above. The remaining sections give the complete scalar
logic and arithmetic bridge. Their separate files are also in the archive.
The local files \file{successor_seed_prelogic.tex} and
\file{successor_seed_bootstrap.tex} are exact copies of
\file{Research/successor_seed_completion/theory/prelogic.tex} and
\file{Research/successor_seed_completion/theory/bootstrap.tex}, respectively.


% Native bridge audit and the pre-selector finite-case repair.
% Requires amsmath, amssymb, amsthm and theorem/lemma environments.

\section{Native bootstrap: cases before a full selector}
\label{nb:section}

This section checks a dependency that matters for the new scalar basis.
Before predecessor is available, the partial selector
\[
 B(n,a,b)=J(n,a;u.b),\qquad B(0,a,b)=Sa,\quad B(Sn,a,b)=b
\]
is not the unrestricted selector $C(0,a,b)=a$, $C(Sn,a,b)=b$.
Accordingly, one must not silently invoke the later ordinary-connective
macros during the arithmetic bootstrap.  In particular, the theorem rule
for numerical cases is not by itself a rule for eliminating equations
$t=0$ and $t=Sz$ from a retained context.  The finite-case lemma below
supplies exactly the stronger form used by the bootstrap.

We use the native Eq/$Q$/$D$ presentation with hypothesis use, full
formula transport, $Q$ elimination, paired elimination, and the two
paired introductions.  In notation,
\[
 Q_yF=\text{the clause taking a proof of }F(0)\text{ to }F(t),
\]
\[
 \frac{\Gamma,F(t,0)\vdash F(t,Sz)}
      {\Gamma\vdash D_{x,y}F(x,y)},\qquad
 \frac{\Gamma,E(t,z)\vdash E(t,Sz)}
      {\Gamma\vdash D_{x,y}E(x,y)}.
\]
The second matrix is a literal equation.  The usual freshness conditions
are retained, and paired elimination opens $Q_yF(a,y)$ at a fresh $a$.
Computation equations may be read as leaves, or obtained from their
constant generated-assumption instances and hypothesis use.

\subsection{Resources available without a numerical selector}

Choose $v$ fresh for $t$ and put $r=J(S0,0;v.t)$.  Computation gives
$r=t$.  Transport this equation along itself with motive $w=t$ to obtain
$t=t$.  The familiar equality motives then give symmetry, transitivity
and numerical congruence; this uses neither predecessor nor successor
injectivity.

For any native formula template $F$, case introduction with an unused
witness coordinate gives $D_{x,y}F(y)$ from
$\Gamma\vdash F(Sz)$, ignoring the retained $F(0)$.  Paired elimination
opens $Q_yF(y)$; a supplied proof of $F(0)$ followed by $Q$ elimination
at $t$ gives
\[
 \frac{\Gamma\vdash F(0)\qquad\Gamma\vdash F(Sz)}
      {\Gamma\vdash F(t)}. \tag{Native NC}\label{nb:nc}
\]
The identical construction with progress introduction, using its local
$E(z)$, gives contextual literal-equation induction from a base and step.
These arguments use no arithmetic selector or ordinary logical connective.

For reference, $Q_yF$ and $D_{x,y}F$, with $x$ unused, are interderivable
already: use case introduction and $Q$ elimination in the forward
direction, and paired elimination followed by hypothesis use in the
reverse direction.  Thus the additional $Q$ constructor presents the
same internal clauses as the Eq/$D$ formulation.

Even equation explosion is available at this stage.  Under $0=1$,
congruence gives $B(0,0,t)=B(1,0,t)$; computation gives $1=t$.
Consequently every equation follows.  External induction on fixed native
formula syntax, using case introduction (and the unused-witness
representation for $Q$), gives explosion for every native formula.
This is an external proof transformation, not compound-formula induction
inside the arithmetic theory.

\subsection{Functional finite ranges do not assume branch equations}

Say that a particular term $b$ has the \emph{functional range rule}
$\operatorname{FR}_R(b)$, for a finite list of numerals $R$, when the
following metarule is derived for every native formula template $K$:
\[
 \bigl(\Gamma\vdash K(r)\text{ for each }r\in R\bigr)
 \quad\Longrightarrow\quad\Gamma\vdash K(b).
 \tag{FR}\label{nb:fr}
\]
There are no added equations $b=r$ in these premises.  The retained
context and every fixed parameter of $K$ remain unchanged.  In
particular, they may mention the original term $b$.

Use the fixed terms $Z,N,j,M$ defined in
Section~\ref{js:cycles}, with $Z(t)=B(t,0,0)$. The present proof
checks their logical prerequisites; their finite computation identities
are verified directly from those definitions.

For $Z(t)$, apply Native NC to the fresh template $w\mapsto K(Z(w))$.
The zero branch is $K(1)$ and every successor branch is $K(0)$, by
computation and transport.  Hence $\operatorname{FR}_{\{0,1\}}(Z(t))$
holds, and therefore so does the same rule for $N(t)$.
For $B(t,1,0)$ the identical construction has endpoints two and zero.
For
\[
 j(t)=B\bigl(t,0,B(Z(N(t)),1,0)\bigr)
\]
use NC on the outer control: the zero branch is one, and the successor
branch has the preceding two-value rule.  This proves
$\operatorname{FR}_{\{0,1,2\}}(j(t))$, and hence the same rule for
$M(t)$, since $M$ is an outer application of $j$.

The inner parity iterate $U(t)=J(t,0;u.Z(u))$ has a two-value range
rule as well.  At counter zero it is one; at a successor it is a value
of $Z$.  Apply NC on the counter and the already proved range rule of
$Z$.  No induction with a disjunctive invariant is involved.

These constructions justify the cycle equations before contextual
finite-case elimination is used.  For example, $Z(Z(U(t)))=U(t)$ follows
by applying its two-value rule to the equation template $Z(Z(w))=w$.
The equations
\[
 M(0)=0,\quad M(1)=1,\quad M(2)=2,\qquad M(Sw)=j(M(w))
\]
and $j(0)=1$, $j(1)=2$, $j(2)=0$ then follow from the fixed definitions.
Their functional range proofs are independent of the ramp invariant.

\begin{lemma}[Successor cancellation on the three finite values]
\label{nb:finite-cancellation}
Suppose $b$ and $c$ have functional range rules whose values are contained
in $\{0,1,2\}$.  In any retained context,
\[
                         Sb=Sc\quad\Longrightarrow\quad b=c
\]
is derivable, before global successor injectivity is known.
\end{lemma}
\begin{proof}
Apply the two functional range rules to the equation templates
$M(w)=w$ and $j(j(j(w)))=w$.  Their required numeral instances are fixed
computations.  Thus $M(b)=b$, $M(c)=c$, $j^3(b)=b$ and $j^3(c)=c$
are already theorems in the retained context.
Apply $M$ to $Sb=Sc$ and use $M(Sw)=j(M(w))$ to get $j(b)=j(c)$.
Apply $j$ twice and use the two cycle identities to obtain $b=c$.
Every step is ordinary numerical congruence or transport.  In particular,
no cancellation of an arbitrary successor was inserted.
\end{proof}

\begin{lemma}[Contextual finite-case elimination for equations]
\label{nb:guarded-cases}
Let $b$ have $\operatorname{FR}_R(b)$, where $R$ is contained in
$\{0,1,2\}$, and let $F$ be the literal equation $u=v$.
If, for each $r\in R$, there is a derivation
\[
                         \Gamma,b=r\vdash u=v,
\]
then $\Gamma\vdash u=v$ is derivable.  The context may contain arbitrary
native formulas, including an induction hypothesis that depends on $b$.
\end{lemma}
\begin{proof}
Define a native formula template, with both bound variables fresh,
\[
 K(x):=D_{a,y}\bigl(B(y,x,u)=B(y,b,v)\bigr),
\]
where $a$ is unused.  Fix $r\in R$.  To introduce $K(r)$ by the case
choice with witness zero, its local zero assumption reduces by
computation to $Sr=Sb$.  Lemma~\ref{nb:finite-cancellation}, with the
constant $r$ and the known finite-range term $b$, yields $b=r$.
Cut in the given branch derivation.  Its conclusion $u=v$ is exactly
the successor matrix, up to the computations of $B$.  Thus
$\Gamma\vdash K(r)$.

Apply the functional range rule of $b$ to these closed branch proofs.
It gives $\Gamma\vdash K(b)$, whose matrix is
$B(y,b,u)=B(y,b,v)$.  Paired elimination opens its $Q$ clause.
At zero its two sides are $Sb$, so reflexivity and computation supply
the base.  Specializing the clause at one gives $u=v$.  Close paired
elimination.  All temporary coordinates are fresh and all temporary
assumptions have been discharged.
\end{proof}

The derivation is a finite guard construction, not a primitive implication
connective.  It relies only on the partial selector, computations,
transport, hypotheses, case introduction, paired elimination and $Q$
elimination.  Neither predecessor nor the unrestricted selector occurs.

\subsection{Audit of the bootstrap and completion of the bridge}

The dependency order is now explicit.  First establish functional range
rules and the two finite cycles using only Native NC and equations.
Then establish finite successor cancellation and
Lemma~\ref{nb:guarded-cases}.  The finite splits in the ramp induction,
the affine scaling induction and the factor-map calculation are instances
of that lemma, with $b=N(n)$ or another displayed finite-range term and
with a literal-equation conclusion.  Retained induction hypotheses are
therefore legitimate.  The saturation identity $K(c,Su)=S K(c,u)$ is
proved directly by NC with no branch assumptions; it does not require
an unrestricted contextual split on arbitrary $c$.  No step in this order
uses the later ordinary logical macros.

After the corrected predecessor $P$ is constructed, the initializer
compiler gives the ordinary scalar $I$.  Consequently
\[
 C(n,a,b):=I(n,a;u.b)
\]
has the unrestricted selector endpoints.  The inherited gate, padding,
formula-selector and ordinary-connective constructions may now be applied
in their stated order.  Binder congruence itself uses only case
introduction, paired elimination, transport and cut; it is already valid
before this final stage.

For clarity, the ordinary interpretations of the native binders are
\[
 (Q_yB)^\circ=B^\circ(0)\to\forall y\,B^\circ(y),\qquad
 (D_{x,y}B)^\circ=
 \exists x\bigl(B^\circ(x,0)\to\forall y\,B^\circ(x,y)\bigr).
\]
Once the logical macros are available, the native return translations
are proved by the same explicit constructions:
\[
 Q_yB\dashv\vdash B(0)\Rightarrow_*\forall_*y\,B(y),
\]
\[
 D_{x,y}B\dashv\vdash
 \exists_*x\bigl(B(x,0)\Rightarrow_*\forall_*y\,B(x,y)\bigr).
\]
Forwards, an opened $Q$ clause under its base supplies the arbitrary
instance needed by universal introduction.  Backwards, a displayed
implication and universal supply the successor matrix for case
introduction; an unused witness gives the $Q$ case.  The existential
witness is retained in the $D$ case.  External structural induction on
finite formulas, together with the endpoint proofs for the ordinary
macros, yields both formula composites.  Falsity uses the explicit
ordinary separation bridge $(0=1)\to\bot$.

Finally, the inherited scalar completion and fixed-description PR
compiler are used only after the logical interface and initialized $I$
have been recovered.  Recurrence uniqueness gives the numerical
composites on the whole original domain.  The domain is unchanged,
interpreted equality is actual equality, and the numerical isomorphism
graph is $x=y$.  Thus the finite-case repair closes the pre-selector
logical dependency; it does not replace the separately supplied complete
arithmetic and logical interpretations by standard-model calculations.


% Successor-seeded scalar iteration: a self-contained arithmetic bootstrap.
% Requires amsmath, amssymb, amsthm and theorem/lemma/proposition environments.
% No document wrapper. All induction invariants below are literal equations.

\section{A successor-seeded scalar basis with two computation schemas}
\label{js:section}

This section completes the explicitly unresolved one-variable
successor-seeded candidate in the supplied shared-continuation
report~\cite{shared}, in its section ``A genuinely smaller numerical
grammar: the precise unresolved step.''  The candidate basis and its
conditional initializer compiler are inherited from that report.
Unlike the counter-aware successor recursor, its update has just one
changing input.  The initializer remains a separate argument, but the
iterator starts at its successor.  The missing result is a numerical
predecessor term with a whole-domain proof, obtained here by adapting
the supplied scalar $2/3$ factor-map construction to shifted arithmetic
and correcting a saturated predecessor at three small arguments.

\subsection{Syntax, computations, and the precise resources used}

The terms are finite staged expressions
\[
 t ::= x\mid 0\mid J(n,a;u.h).
\]
Only $h$ binds $u$; its other free variables are passive parameters.
Define the abbreviation
\[
 Sa:=J(0,a;u.0),\qquad 1:=S0,
\]
and use standard numerals as finite strings of this abbreviation.
There are exactly two numerical computation schemas:
\begin{align}
 J(0,a;u.h)&=Sa,\tag{J0}\label{js:j0}\\
 J(Sn,a;u.h)&=h[J(n,a;u.h)/u].\tag{JS}\label{js:js}
\end{align}
The first schema, with the abbreviation expanded, says that
$J(0,a;u.h)=J(0,a;u.0)$.  It does not presuppose a separate successor
constructor.  No function variables or uniform evaluator are present.

The bootstrap uses equality replacement, numerical cases
\[
 \frac{\Gamma\vdash F(0)\qquad\Gamma\vdash F(Sz)}
      {\Gamma\vdash F(t)},
 \tag{NC}\label{js:nc}
\]
and contextual induction only on a literal equation $E$:
\[
 \frac{\Gamma\vdash E(0)\qquad\Gamma,E(z)\vdash E(Sz)}
      {\Gamma\vdash E(t)}.
 \tag{EI}\label{js:ei}
\]
The ordinary eigenvariable conditions apply.  In an ordinary
intuitionistic first-order presentation, NC follows from coverage
$\forall x\,(x=0\vee\exists y\,x=Sy)$.  In the native paired-binder
presentation it is the already available case rule.  No use of
successor injectivity, predecessor, pairing, an equality characteristic,
or induction on a compound formula is made before predecessor is
constructed below.

The finite range splits used below require the distinction proved in
Section~\ref{nb:section}. Functional range elimination follows from NC;
contextual finite cases with retained induction hypotheses use the finite
cancellation and guard construction of Lemma~\ref{nb:guarded-cases}.
The cycle identities are established first by functional range elimination
alone. Thus a displayed finite range is not an uncharged disjunction or
an assumption of the ordinary-connective macros.

We shall use the following admissible consequence of EI.  If two fixed
updates are pointwise equal, and the displayed initial values agree,
their iterates agree at every counter.  Induct on the literal equation
between the two iterates: J0 supplies its base and JS supplies its step.
This is a derived uniqueness argument, not a primitive extensionality
rule for update bodies.  All passive arguments are kept fixed while
performing the induction.

\subsection{A partial selector and shifted addition}

Put
\[
 B(n,a,b):=J(n,a;u.b),\qquad Z(n):=B(n,0,0),
 \qquad a\oplus b:=J(b,a;u.Su).
\]
Their computation endpoints are
\[
 B(0,a,b)=Sa,\quad B(Sn,a,b)=b,\qquad
 Z(0)=1,\quad Z(Sn)=0,
\]
\[
 a\oplus0=Sa,\qquad a\oplus Sb=S(a\oplus b).
\]
The operation $\oplus$ has standard meaning $a+b+1$.  Ordinary
unshifted addition has not been defined and is not used in the proof.
Literal-equation induction proves
\[
 0\oplus a=Sa,\quad Sa\oplus b=S(a\oplus b),\quad
 a\oplus b=b\oplus a,\quad
 (a\oplus b)\oplus c=a\oplus(b\oplus c).
 \tag{Shifted algebra}\label{js:oplus}
\]
For the first two equations induct on the second argument.  For
commutativity induct on $b$, using the first two equations.  For
associativity induct on $c$; at zero both sides are
$S(a\oplus b)$, and the successor step applies $S$ to the previous
equality.  These proofs do not cancel a successor.

In particular, the following are genuine available terms:
\[
 D(a):=S^2(a\oplus a),\qquad
 T(a):=S^4((a\oplus a)\oplus a),\qquad
 R(a):=S^3(a\oplus a).
 \tag{Affine terms}\label{js:affine}
\]
Their standard meanings are respectively $2a+3$, $3a+6$, and
$2a+4$.  The proof uses their displayed term definitions.  From
\eqref{js:oplus} one obtains
\[
 D(0)=3,\quad D(Sa)=S^2D(a),\qquad
 T(0)=6,\quad T(Sa)=S^3T(a).
 \tag{Affine shifts}\label{js:affine-shifts}
\]

\subsection{Two finite cycles}\label{js:cycles}

Define
\[
 N(n):=Z\bigl(J(n,0;u.Z(u))\bigr).
\]
The inner iterator is always zero or one: at zero it is one, and at a
successor it is a value of $Z$.  The endpoints of $Z$ and finite cases
therefore give
\[
 N(0)=0,\quad N(Sn)=Z(N(n)),\quad N(S^2n)=N(n),
 \qquad N(n)\in\{0,1\}.
 \tag{Parity cycle}\label{js:parity}
\]
This asserts a proved two-state recurrence, not a pre-existing parity
operation.  For example, if $K(n)=J(n,0;u.Z(u))$, then
$N(Sn)=Z(Z(K(n)))=K(n)=Z(N(n))$ by its two cases.

Next put
\[
 j(u):=B\bigl(u,0,B(Z(N(u)),1,0)\bigr),\qquad
 M(n):=j\bigl(j(J(n,0;u.j(u)))\bigr).
\]
Finite computations give
\[
 j(0)=1,\qquad j(1)=2,\qquad j(2)=0.
\]
The range of $j$ is contained in $\{0,1,2\}$: its zero branch is
one, and its successor branch is zero or two according to $N(u)$.
Consequently the definitions and these three computations give
\[
 M(0)=0,\quad M(Sn)=j(M(n)),\quad M(S^3n)=M(n),
 \qquad M(n)\in\{0,1,2\}.
 \tag{Three-cycle}\label{js:three-cycle}
\]
Here $M(0)=j(j(1))=0$, and the successor equation follows by
associativity of composing the fixed term $j$ with itself.  The period
three follows from its three possible values.  No numerical division
algorithm has been assumed.

\subsection{The shifted ramp and factor map}

Set
\[
 G(n):=J(n,0;u.u\oplus M(S^2u)).
\]
Its first values are
\[
 G(0)=1,\ G(1)=2,\ G(2)=4,\ G(3)=5,\
 G(4)=7,\ G(5)=8,\ G(6)=10.
\]

\begin{lemma}[Ramp identities]\label{js:ramp}
The preceding theory proves
\[
 M(G(n))=S(N(n)),\qquad G(S^2n)=S^3G(n),
 \tag{Ramp}\label{js:ramp-identities}
\]
and admits the following two finite-case calculations:
\[
 \begin{array}{c|cc}
 &N(n)=0&N(n)=1\\ \hline
 G(Sn)&S G(n)&S^2G(n)\\
 G(S^3n)&S^4G(n)&S^5G(n).
 \end{array}
 \tag{Ramp cases}\label{js:ramp-cases}
\]
\end{lemma}
\begin{proof}
Induct on the single equation $M(G(n))=S(N(n))$.
Its base is $M(1)=1$.  If $N(n)=0$, the induction equation gives
$M(G(n))=1$, hence $M(S^2G(n))=0$ by the three-cycle.  The next
state is $G(n)\oplus0=S G(n)$, whose $M$-value is two, as
required.  If $N(n)=1$, the current $M$-value is two and
$M(S^2G(n))=1$.  The next state is
$G(n)\oplus1=S^2G(n)$, whose $M$-value is one.  This proves
the invariant and the first line of the case table.  Applying that line
twice, with the parity cycle, gives the two-step equation; applying it
three times gives the last line.
\end{proof}

\begin{lemma}[Affine identities]\label{js:affine-lemmas}
Throughout the entire numerical domain,
\[
 N(D(a))=1,\qquad N(T(a))=N(a),\qquad
 S G(D(a))=T(a),
 \tag{Affine parity and ramp}\label{js:affine-parity}
\]
\[
 G(T(a))=T(G(a))\oplus N(a).
 \tag{Ramp scaling}\label{js:ramp-scaling}
\]
\end{lemma}
\begin{proof}
For $N(D(a))=1$, use EI, $D(0)=3$, and the two-step parity
period.  For $N(T(a))=N(a)$ use EI, $T(0)=6$, and
$T(Sa)=S^3T(a)$; three parity steps toggle both sides once.

For $S G(D(a))=T(a)$ use EI on that literal equation.  The base
is $S G(3)=6$.  The successor step uses
$D(Sa)=S^2D(a)$ and $G(S^2v)=S^3G(v)$, so it applies three
successors to the previous equation, just as $T$ does.

For \eqref{js:ramp-scaling}, the base is
$G(6)=10=T(1)\oplus0$.  Retain its equation as the induction
hypothesis, and split $N(a)$.  In the zero case, $N(T(a))=0$,
so the new left side is $S^4G(T(a))$.  The old right side is
$T(G(a))\oplus0=S T(G(a))$, while the new right side is
\[
 T(SG(a))\oplus1=S^5T(G(a)).
\]
Thus both sides have changed by four successors.  In the one case,
the new left side is $S^5G(T(a))$.  The old right side is
$T(G(a))\oplus1=S^2T(G(a))$, while the new right side is
\[
 T(S^2G(a))\oplus0=S^7T(G(a)).
\]
Thus both sides have changed by five successors.  These calculations
use only \eqref{js:oplus}, \eqref{js:affine-shifts}, and the ramp
case table; no successor has been cancelled.
\end{proof}

Define the factor map
\[
 H(u):=B(Z(N(u)),G(u),u).
\]
Its two cases are
\[
 N(u)=0\ \Longrightarrow\ H(u)=u,\qquad
 N(u)=1\ \Longrightarrow\ H(u)=S G(u).
\]

\paragraph{Why these shifts are useful.}
For intuition on standard natural numbers only, $G(u)=\lfloor3u/2\rfloor+1$.
Thus $H$ fixes even inputs and sends an odd $u$ to $3(u+1)/2$.
After shifting the state by three, an active step multiplies it by $3/2$.
The affine maps $D(a)=2a+3$ and $T(a)=3a+6$ are precisely doubling
and tripling in this shifted coordinate: $D(a)+3=2(a+3)$ and
$T(a)+3=3(a+3)$. Starting at two makes the unshifted base even and
fixed by $H$, while its shifted value is five. The later seed sequence
therefore has standard values $5\cdot2^x-3$. This motivates a comparator
that tests when its powers of two have been exhausted. The proof below
uses the displayed $J$ terms and literal equations; division, exponentiation,
and this standard-value description are not prerequisites.

\begin{lemma}[The factor-map equations]\label{js:factor}
The following global equations hold:
\[
 H(2)=2,\qquad H(D(a))=T(a),\qquad H(T(a))=T(H(a)).
 \tag{Factor}\label{js:factor-equations}
\]
\end{lemma}
\begin{proof}
The first equation follows from $N(2)=0$.  Since $N(D(a))=1$,
the second is $S G(D(a))=T(a)$.
For the last equation split $N(a)$, noting that $N(T(a))=N(a)$.
In the zero case both sides equal $T(a)$.  In the one case,
\eqref{js:ramp-scaling} gives
\[
 H(T(a))=S G(T(a))
 =S\bigl(T(G(a))\oplus1\bigr)
 =S^3T(G(a))=T(SG(a))=T(H(a)).
\]
\end{proof}

\subsection{A comparator without stripping an initializer}

Put
\[
 A(x):=J(x,0;u.R(u)),\qquad E(x):=S A(x),
\]
\[
 F(x,y):=J(y,A(x);u.H(u)),\qquad
 L(x,y):=Z(N(F(x,y))).
\]
By direct computation and shifted algebra,
\[
 E(0)=2,\qquad E(Sx)=D(E(x)),\qquad F(x,0)=E(x).
\]
The definition of $F$ is the decisive initializer alignment:
its successor-seeded initial value is exactly $E(x)=S A(x)$.
There is no hidden use of predecessor to initialize it at $E(x)$.

\begin{lemma}[Comparator equations]\label{js:comparator}
The term $L$ is zero or one and satisfies
\[
 L(0,y)=1,\qquad L(Sx,0)=0,\qquad
 L(Sx,Sy)=L(x,y).
 \tag{Comparator}\label{js:comparator-equations}
\]
\end{lemma}
\begin{proof}
Its range follows from the range of $Z$.  EI on $y$ proves
$F(0,y)=2$, since $E(0)=2$ and $H(2)=2$; hence $L(0,y)=1$.
Also $F(Sx,0)=E(Sx)=D(E(x))$ has parity one, giving
$L(Sx,0)=0$.

For the remaining equation first prove, by EI on $y$ with $x$ fixed,
\[
 F(Sx,Sy)=T(F(x,y)).
\]
At zero, the left side is
$H(E(Sx))=H(D(E(x)))=T(E(x))$.  In the successor step both
iterates apply $H$, and $H\circ T=T\circ H$ rewrites the result
into $T(F(x,Sy))$.  Applying $Z\circ N$ and using
$N(T(v))=N(v)$ gives the desired comparator equation.

This proof avoids cancellation in an equation between two successor
initializers.  Such cancellation would be unavailable at this stage.
\end{proof}

\subsection{Saturation and a finite correction give true predecessor}

Define conditional increment by
\[
 K(c,u):=B(c,u,u),\qquad
 B_m(u):=K(Z(L(Su,m)),u),\qquad
 V(t,m):=J(t,0;u.B_m(u)).
\]
Thus $K(0,u)=Su$ and $K(Sv,u)=u$.  The counter $m$ is a passive
parameter of the iteration defining $V$.

\begin{lemma}[Shifted saturation]\label{js:saturation}
The theory proves
\[
 V(St,S^2k)=S V(t,Sk),\qquad
 V(n,Sn)=Sn,\qquad
 V(Sn,Sn)=V(S^2n,Sn)=Sn.
 \tag{Saturation}\label{js:saturation-equations}
\]
\end{lemma}
\begin{proof}
Numerical cases on $c$ give $K(c,Su)=S K(c,u)$.  Consequently
the comparator equations give
\[
 B_{Sm}(Su)=S B_m(u),\qquad B_{S^2k}(1)=2.
\]
EI on $t$, with $k$ fixed, now proves the first displayed equation.
At zero its two sides equal two.  At a successor, apply the defining
update to the induction equation and then use the conjugacy
$B_{S^2k}(Su)=S B_{Sk}(u)$.

EI gives $L(Sn,n)=0$ directly from the comparator equations.
A separate EI proves $V(n,Sn)=Sn$: the base is $V(0,1)=1$,
and the successor step uses the previously proved uniform shift at
$t=k=n$.  This is substitution into a completed lemma, not a change
of a passive parameter inside an active induction hypothesis.
Finally
\[
 B_{Sn}(Sn)=K(Z(L(S^2n,Sn)),Sn)
           =K(Z(L(Sn,n)),Sn)=K(1,Sn)=Sn.
\]
One and then two more iterator steps give the last two equations.
\end{proof}

Put
\[
 Q(m):=J\bigl(m,0;u.K(Z(L(S^2u,m)),u)\bigr).
\]
This term is a capped predecessor rather than the desired predecessor:
\[
 Q(0)=Q(1)=1,\qquad Q(S^2k)=Sk.
 \tag{Capped predecessor}\label{js:capped}
\]
The first two equations are direct computations.  For the last,
\[
 L(S^2u,S^2k)=L(u,k)=L(Su,Sk)
\]
identifies its update with $B_{Sk}(u)$.  The derived iterator uniqueness
lemma, proved at an independent counter $t$ with $k$ fixed, gives
\[
 J(t,0;u.K(Z(L(S^2u,S^2k)),u))=V(t,Sk).
\]
Substitute $t=S^2k$ and apply \eqref{js:saturation-equations}.

Now define one fixed correction term
\[
 \boxed{\quad
 P(m):=B\bigl(Z(L(3,m)),\ Q(Q(m)),\ L(2,m)\bigr).
 \quad}\tag{Predecessor}\label{js:predecessor}
\]
There is no arbitrary selector hidden in this expression.  Its outer
operation is the original partial selector $B$, whose zero branch adds
exactly one successor.

\begin{theorem}[Predecessor and initializer recovery]\label{js:recovery}
The two-schema successor-seeded scalar theory proves
\[
 P(0)=0,\qquad P(Sn)=n.
\]
For every fixed update $h$, the finite term
\[
 I(n,a;v.h):=
 P\bigl(J(n,a;u.S(h[Pu/v]))\bigr)
 \tag{Initializer compiler}\label{js:initializer}
\]
satisfies
\[
 I(0,a;v.h)=a,\qquad I(Sn,a;v.h)=h[I(n,a;v.h)/v].
\]
\end{theorem}
\begin{proof}
The comparator equations give the three direct values
\[
 P(0)=0,\qquad P(1)=0,\qquad P(2)=1.
\]
Indeed, at these inputs $L(3,m)=0$, so the outer $B$ selects
$L(2,m)$, whose three values are zero, zero, one.
At $m=S^3k$, one has $L(3,m)=1$.  Thus the zero-control branch
of $B$, together with \eqref{js:capped}, gives
\[
 P(S^3k)=S Q(Q(S^3k))=S Q(S^2k)=S^2k.
\]
Numerical cases twice on $n$ now give $P(Sn)=n$ from these displayed
equations.  Alternatively, two literal-equation inductions, first on
$P(S^2n)=Sn$ and then on $P(Sn)=n$, give the same conclusion.
No earlier cancellation or injectivity was used.  Successor injectivity
can now be derived by applying $P$ to an equality of successors.

For the initializer compiler, J0 and the proved predecessor equation
give the base $P(Sa)=a$.  At a successor, JS exposes the term
$S(h[P(J(n,a;u.S(h[Pu/v])))/v])$.  The outer $P$ removes its
displayed successor and leaves exactly $h[I(n,a;v.h)/v]$.
There is no invariant about the internal state to prove: the displayed
successor in the update is sufficient.
\end{proof}

\subsection{Consequences and the bi-interpretation boundary}

The construction supplies the full initialized scalar basis $0,S,I$.
In particular the unrestricted selector is now the derived term
$C(n,a,b):=I(n,a;u.b)$, with its ordinary endpoints $a$ and $b$.
The supplied scalar arithmetic development can therefore be used from
this point onward to obtain actual-equality reflection, induction on
arbitrary quantifier-free formulas, pairing and projections, and the
compiler for every fixed general primitive-recursion description.
Those later arguments are not premises of this bootstrap.

In ordinary iPRA, each fixed $J$ description is interpreted by primitive
recursion with seed $Sa$ and its displayed scalar update.  Both
computation schemas, coverage, separation, and literal-equation induction
are then provable.  Conversely \eqref{js:initializer}, followed by the
existing scalar compiler, interprets every fixed iPRA operation.
For each fixed description, literal-equation induction proves uniqueness
of its defining recurrence after earlier subdescriptions have been
identified.  Applied to both composite translations, this identifies
each numerical operation with its double translation throughout the
entire original domain.  Both domain formulas are $x=x$ and both
numerical isomorphism graphs are $x=y$.

For the ordinary presentations, full intuitionistic logic is preserved.
For the native paired-binder presentation, its ordinary logical macros
and the formula return-translation proof must still be included; they
apply once the derived selector $C$ is available.  This establishes
bi-interpretability only when that logical bridge and the scalar
arithmetic completion accompany the present argument.  It does not
replace them by agreement merely on the standard natural numbers.

The arithmetic advance is completion of the predecessor obligation
left open for this exact candidate in the shared-continuation
report~\cite{shared}: the affine factor-map alignment, capped
predecessor, and finite correction \eqref{js:predecessor} supply its
missing proof.  The candidate grammar and initializer compiler are not
new.  Finite cycles, a $2/3$ factor map, and saturation are adapted from
the supplied ordinary scalar predecessor construction.  No global
literature priority is asserted.

The now-completed basis has two primitive numerical constructors and
two computation schemas, with a single update variable.  A smaller
complete proof checker remains a separate implementation claim, which
must charge the actual decoder, scope checks, logical rules, and driver.


% New mathematical result: complementary introduction choices.
% Requires amsmath/amssymb/amsthm and theorem, lemma, proposition environments.
% The inherited arithmetic-to-iPRA bridge must be supplied separately.

\section{Complementary introductions with one local assumption}
\label{sec:choice-introduction}

The old paired introduction retains both the zero instance and an
induction instance when its matrix is an equation. Both together are
unnecessary. A certificate may choose either one, provided the induction
choice is restricted to a literal equation. This changes neither the
formula language nor the meaning of a formula. The omitted combination
is recovered by a short nested introduction--elimination construction.

\subsection{Complete native syntax and rule inventory}

The scalar numerical language is
\[
 t::=x\mid0\mid St\mid I(n,a;u.h).
\]
Only $h$ binds the state variable $u$; $n$ and $a$ do not.
The update may contain fixed numerical parameters. The computation
axioms in every retained context are
\[
 I(0,a;u.h)=a,\qquad
 I(Sn,a;u.h)=h[I(n,a;u.h)/u].
 \tag{I0, IS}\label{eq:choice-computations}
\]
The argument also applies unchanged to the full fixed-description PR
language with its projection, composition, and recursion equations.
No numerical variable ranges over function descriptions, and no uniform
evaluator is postulated.

Native formulas and reusable context clauses are
\[
 A::=(s=t)\mid\mathsf D_{x,y}A,\qquad
 \langle A;y.B(y)\rangle,\quad y\notin\operatorname{FV}(A).
\]
Bound variables are renamed before capture-free substitution.
The two variables $x,y$ in a paired formula are distinct.
An ordinary assumption $F$ is the constant clause
$\langle H;v.F\rangle$, where $v$ is absent from $F$ and $H$ is the
fixed computation theorem $I(0,0;u.0)=0$. Thus ordinary hypothesis
use is a clause Use preceded by this computation proof.

There are computation axioms and exactly the following inference
patterns:
\[
\frac{\Gamma\vdash A}{\Gamma\vdash B(t)}
\quad\mathsf{Use}
\quad(\langle A;y.B(y)\rangle\in\Gamma),
\]
\[
\frac{\Gamma\vdash\mathsf D_{x,y}B(x,y)\qquad
      \Gamma,\langle B(a,0);y.B(a,y)\rangle\vdash F}
     {\Gamma\vdash F}
\quad\mathsf{DE},
\qquad
\frac{\Gamma\vdash s=t\qquad\Gamma\vdash F(s)}
     {\Gamma\vdash F(t)}
\quad\mathsf{Tr},
\]
\[
\frac{\Gamma,B(t,\eta_\varepsilon(z))\vdash B(t,Sz)}
     {\Gamma\vdash\mathsf D_{x,y}B(x,y)}
\quad\mathsf{DI}_{\varepsilon},
\qquad
\eta_0(z)=0,\quad\eta_1(z)=z.
\tag{Choice DI}\label{eq:choice-di}
\]
Here $\varepsilon$ is a certificate choice in $\{0,1\}$, and
$\varepsilon=1$ is permitted only when $B$ has a literal equation
root. An equation beneath a binder is not a literal equation.
Call $\varepsilon=0$ the \emph{case choice} and $\varepsilon=1$
the \emph{progress choice}. Both retain exactly one ordinary
formula assumption. Compound matrices have only the case choice.

In DI the eigenparameter $z$ is fresh for $\Gamma$, the witness
$t$, and every fixed parameter of $B$. In DE the eigenparameter
$a$ is fresh for $\Gamma$, the conclusion $F$, and every fixed
parameter of $B$. In transport, $F(v)$ is any fixed native
formula template. Displayed local assumptions and clauses are
discharged. There are no primitive reflexivity, coverage,
separation, ordinary logical, or additional induction rules.

We call this calculus $\mathcal C$. The original comparison
calculus $\mathcal N$ has identical syntax, computation,
Use, DE, and transport, but its introduction is
\[
\frac{\Gamma,B(t,0),[B(t,z)]_{\mathrm{at}}\vdash B(t,Sz)}
     {\Gamma\vdash\mathsf D_{x,y}B(x,y)}.
\tag{Old DI}\label{eq:choice-old-di}
\]
The bracket supplies $B(t,z)$ exactly at a literal equation.

\subsection{Exact equivalence and the unused-witness construction}

\begin{lemma}[Structural admissibility]
\label{lem:choice-structural}
Weakening, exchange, contraction, renaming, ordinary-assumption cut,
and capture-free numerical substitution are admissible in $\mathcal C$.
\end{lemma}
\begin{proof}
Induct on derivations. Rename local eigenparameters before inserting
proofs or terms. Numerical substitution preserves a formula's root
constructor, so an allowed progress choice remains allowed.
The selected choice does not change under substitution or cut.
For ordinary-assumption cut, replace each use of its constant clause
by the supplied proof. Computation schemas are closed under
capture-free numerical substitution. These facts verify every case.
\end{proof}

\begin{theorem}[Identity-language equivalence]
\label{thm:choice-equivalence}
For every retained clause context $\Gamma$ and native formula $F$,
\[
 \Gamma\vdash_{\mathcal C}F
 \quad\Longleftrightarrow\quad
 \Gamma\vdash_{\mathcal N}F.
\]
The transformations are effective on finite derivations and leave
all numerical and formula syntax unchanged.
\end{theorem}
\begin{proof}
Each new introduction is an instance of the old one after weakening
by the omitted local assumption. For the case choice at an equation,
add the induction instance. For the progress choice, add the zero
instance. Compound case introduction is already the old rule.
Thus every new derivation translates to an old one.

For the reverse direction, transform an old derivation recursively.
Only its atomic DI case needs a construction. Suppose its matrix is
the equation $E(x,y)$ and its transformed premise is
\[
 \pi:\quad\Gamma,E(t,0),E(t,z)\vdash E(t,Sz).
\]
The outer new inference will be case DI with witness $t$. It is
therefore enough to prove $E(t,Sz)$ under
$\Gamma,E(t,0)$.

Rename the eigenvariable in $\pi$ to a fresh $v$. Choose a new
variable $u$ absent from $E(t,y)$ and all retained data. Under
$\Gamma,E(t,0)$, progress DI with witness zero introduces
\[
             \mathsf D_{u,y}E(t,y)
\]
from exactly the renamed premise $\pi[v/z]$.
The bound coordinate $u$ is deliberately unused.

Apply DE to this paired formula. Because $u$ is unused, the opened
clause is exactly
\[
               \langle E(t,0);y.E(t,y)\rangle;
\]
it contains no new existential witness. The retained ordinary
assumption $E(t,0)$ supplies its antecedent. Use at $Sz$ gives
$E(t,Sz)$, and DE closes. Finally the outer case DI introduces
the original $\mathsf D_{x,y}E(x,y)$ with witness $t$.
The original premise proof $\pi$ has been inserted once.

An old compound introduction already is the new case choice.
All other rule cases are unchanged. This completes the induction
on the original derivation.
\end{proof}

The unused inner coordinate is essential. Introducing
$\mathsf D_{x,y}E(x,y)$ internally and then eliminating it would
give an arbitrary fresh witness, not necessarily the original
$t$. The construction instead keeps $t$ as a fixed parameter
of $E(t,y)$ and binds only the unused coordinate $u$.

\begin{proposition}[Linear bound on inference nodes]
\label{prop:choice-node-bound}
Count each computation axiom, Use, DI, DE, and transport as one
inference node, expanding ordinary-assumption use as a computation
of $H$ followed by Use. If an old derivation has $N$ nodes and
$r$ atomic DI nodes, the displayed transformation has exactly
\[
                   N+5r\leq 6N
\]
nodes.
\end{proposition}
\begin{proof}
Each atomic old DI node is replaced by six nodes: outer case DI,
inner progress DI, DE, Use of the opened clause, Use of the
retained base assumption, and the computation of $H$ needed for
that last ordinary-assumption use. Its old premise subtree
occurs once. All other inference nodes are copied once.
\end{proof}

This counts inference-tree nodes. It is not a bound on serialized
JSON bytes, formula sizes, checker runtime, or the number of
consumed packet cells when hypothesis closure is implicit.
In the formula-context Q/D presentation, the corresponding
replacement has outer case DI, inner progress DI, DE, QE,
and two Hyp nodes, giving the same six-node count.

\subsection{Induction and numerical cases separate cleanly}

\begin{lemma}[Contextual equation induction]
\label{lem:choice-ei}
The progress choice, DE, and Use derive
\[
\frac{\Gamma\vdash E(0)\qquad
      \Gamma,E(z)\vdash E(Sz)}
     {\Gamma\vdash E(t)}
\tag{EI}\label{eq:choice-ei}
\]
for every equation $E$, with the usual freshness condition.
\end{lemma}
\begin{proof}
Choose an unused witness coordinate $u$. Progress DI introduces
$\mathsf D_{u,y}E(y)$ from the given step. DE opens the clause
$\langle E(0);y.E(y)\rangle$. Supply its antecedent using the given
base proof and use the clause at $t$. Close DE.
\end{proof}

\begin{lemma}[Cases, equality, and binder congruence]
\label{lem:choice-preinduction}
The case choice derives, for arbitrary native formulas,
\[
\frac{\Gamma\vdash F(0)\qquad\Gamma\vdash F(Sz)}
     {\Gamma\vdash F(t)}
\tag{NC}\label{eq:choice-nc}
\]
with the usual freshness condition. Equality reflexivity,
symmetry, transitivity, numerical congruence, and uniform
matrix interderivability under $\mathsf D$ are derivable
without using the progress choice.
\end{lemma}
\begin{proof}
For cases, use case DI with an unused witness coordinate and
matrix $F(y)$, ignoring its retained base in the successor
premise. DE opens $\langle F(0);y.F(y)\rangle$; the given zero
proof and Use at $t$ finish.

For reflexivity put $r=I(0,t;u.0)$. Use the computation equation
$r=t$ as both transport premises with motive $v=t$ to obtain
$t=t$. The usual explicit equality motives give symmetry,
transitivity, and numerical congruence.

For binder congruence, open the old paired formula at a fresh
witness $a$ and introduce the new one by case DI at $a$.
Convert its retained base to the old base, use the opened
old clause at the successor, and convert back. Ordinary cut
closes the construction, and the reverse conversion gives
the converse. No induction is used.
\end{proof}

The primitive progress choice is still restricted to equations.
Allowing it at arbitrary compound matrices would, by the
argument for \eqref{eq:choice-ei}, supply induction for
arbitrary formulas and change the intended theory.

\subsection{The entire ordinary logical interface}
\label{subsec:choice-logic}

Here is a predecessor-free reconstruction of the ordinary logical
operations. It makes explicit that the case choice supplies the
logical interface independently of progress induction.

Let $C(v,a,b):=I(v,a;u.b)$ with $u$ unused. Computation gives
$C(0,a,b)=a$ and $C(Sv,a,b)=b$. Define the gate on fixed finite
formula syntax by
\[
 G_v(s=t):=(C(v,0,s)=C(v,0,t)),\qquad
 G_v(\mathsf D_{x,y}B):=\mathsf D_{x,y}G_vB,
\]
renaming bound variables away from the control. Equality,
binder congruence, and external structural induction give
\[
 \vdash G_0A,\qquad G_{Sv}A\dashv\vdash A,\qquad
 A\dashv\vdash\mathsf D_{u,y}G_yA
 \quad(u,y\notin\operatorname{FV}(A)).
\tag{Gate and padding}\label{eq:choice-gate}
\]
For padding, case DI proves the forward direction. DE, the
zero-gate theorem, and Use at $1$ prove the reverse direction.

Define a selector between two arbitrary formulas by
\[
 M_v(s=t,u=w):=(C(v,s,u)=C(v,t,w))
\]
and, with aligned fresh paired variables, by
\[
\begin{aligned}
 M_v(\mathsf D_{x,y}A,\mathsf D_{x,y}B)
   &:=\mathsf D_{x,y}M_v(A,B),\\
 M_v(E,\mathsf D_{x,y}B)
   &:=\mathsf D_{x,y}M_v(G_yE,B),\\
 M_v(\mathsf D_{x,y}A,E)
   &:=\mathsf D_{x,y}M_v(A,G_yE),
\end{aligned}
\]
where $E$ is an equation. Every call decreases maximum input
binder depth. Padding and congruence prove
$M_0(A,B)\dashv\vdash A$ and
$M_{Sv}(A,B)\dashv\vdash B$.

Put $\delta(v)=C(v,1,0)$ and $p(z)=I(z,0;u.\delta(u))$.
Only the fixed computations $p(1)=1$ and $p(2)=0$ are needed.
The full logical macro inventory is
\[
\begin{aligned}
 A\Rightarrow_*B&:=\mathsf D_{u,z}M_z(A,B),\\
 A\wedge_*B&:=\mathsf D_{u,z}G_zM_{p(z)}(A,B),\\
 \exists_*x\,A(x)&:=\mathsf D_{x,y}G_yA(x),\\
 A\vee_*B&:=\exists_*x\,M_x(A,B),\\
 \forall_*x\,A(x)&:=
        A(0)\wedge_*\mathsf D_{u,z}G_zA(z),\\
 \bot_*&:=(0=1).
\end{aligned}
\tag{Logical macros}\label{eq:choice-macros}
\]
Newly bound variables are fresh; $u$ is unused.

Implication introduction obtains its assumption from the zero
selector and supplies the consequent at successors, then uses
case DI. Elimination opens its clause and uses it at $1$.
Existential introduction uses its witness and the successor
gate; elimination opens at a fresh witness, supplies the
zero-gate theorem, and uses the clause at $1$.

Disjunction introduction uses witnesses $0$ and $1$.
For elimination, discharge the two branch proofs as
$A\Rightarrow_*F$ and $B\Rightarrow_*F$. NC applied to
$M_z(A,B)\Rightarrow_*F$ gives this implication at every
control: the two endpoint cases follow from the branch
implications. Existential elimination provides $M_a(A,B)$
and the implication gives $F$.

Given $A$ and $B$, selector endpoints and NC give $M_v(A,B)$
at an arbitrary control. Substitute $v=p(Sz)$ and use the
successor gate for conjunction introduction. For projections,
open its zero-gated clause and specialize at $1$ and $2$;
the two fixed computations give $B$ and $A$, respectively.

These rules prove existential numerical coverage:
\[
     (x=0)\vee_*\exists_*v\,(x=Sv).
\tag{Coverage}\label{eq:choice-coverage}
\]
Use NC. Zero takes the left injection of reflexivity, and a
successor takes the right injection with the evident witness.
Coverage gives equality cases even when the retained context
contains $x$; its elimination adds $x=0$ or $x=Sv$ rather
than silently changing that retained context.

Universal introduction follows from a fresh-parameter proof
by substituting zero and successors and forming the two
conjuncts. For elimination, take its $A(0)$ conjunct and
open the paired conjunct. Its base is a zero-gate theorem.
Coverage of the target chooses the zero instance from
$A(0)$ or a successor instance from the opened clause
and gate removal.

Finally, under $0=1$, computation and transport give
$s=C(0,s,t)=C(1,s,t)=t$ for arbitrary numerical terms.
Every equation follows, and external induction on formula
syntax with case DI gives every native formula. This
supplies falsity elimination. No excluded middle for
quantified formulas has entered the construction.

\subsection{Whole-domain realizations and both formula composites}

For a standard first-order reading, let $T_{\rm eq}$ have
intuitionistic equality logic, the same numerical signature
and computation equations, numerical coverage, separation
$(0=1)\to\bot$, and contextual equation induction.
Its interpretation of the native syntax is
\[
 (s=t)^\circ=(s=t),\qquad
 (\mathsf D_{x,y}B)^\circ
   =\exists x\bigl(B^\circ(x,0)\to\forall y B^\circ(x,y)\bigr),
\]
\[
 \langle A;y.B\rangle^\circ
       =A^\circ\to\forall y B^\circ.
\]
The domain is the entire numerical domain and interpreted
equality is actual equality.

Case DI is sound in this realization by numerical coverage.
Progress DI is sound by contextual equation induction under
the conditional base. Use, DE, transport, and computation
have their displayed ordinary meanings. Conversely, the
native macros derive every ordinary logical rule,
\eqref{eq:choice-ei} derives equation induction, and
\eqref{eq:choice-coverage} derives coverage. Separation
translates to the identity implication
$(0=1)\Rightarrow_*(0=1)$.

For the essential paired formula composite,
\[
 \mathsf D_{x,y}B
 \dashv\vdash
 \exists_*x\bigl(B(x,0)\Rightarrow_*\forall_*y\,B(x,y)\bigr).
\tag{Paired composite}\label{eq:choice-paired-composite}
\]
Forwards, DE opens its clause at a fresh $a$; under $B(a,0)$
that clause supplies $B(a,y)$ at a fresh arbitrary $y$.
Introduce the universal, implication, and existential
macros and close DE. Backwards, eliminate the existential
at a fresh $a$. Introduce the paired formula with witness
$a$ by case DI: the retained base, supplied implication,
and universal elimination at $Sz$ give its successor premise.

The ordinary interpretations of the six logical macros
have their ordinary round trips. For conjunction, its true
base and the instances $1,2$ recover both operands;
the converse uses selector cases. The universal has all
successor instances in its paired conjunct and zero in
its first conjunct; coverage gives exactly the ordinary
universal. Implication, disjunction, and existential
quantification use their selector or gate endpoints.
Falsity uses the explicit ordinary separation bridge.
External induction on finite formulas, together with
\eqref{eq:choice-paired-composite}, gives both formula
composites. Clauses correspond in both directions to
$A\Rightarrow_*\forall_*yB$ by the derived logical rules.

The numerical composites are literal identities with
internally provable graph $x=y$. No standard-numeral
restriction, retract, quotient, or equality reflection
assumption is being inserted at this stage. The separately
supplied arithmetic bridge identifies $T_{\rm eq}$ with
full iPRA; the present change preserves that bridge
unchanged because Theorem~\ref{thm:choice-equivalence}
preserves every native sequent.

\subsection{What the theorem does and does not claim}

The new inference interface retains one ordinary assumption
at every paired introduction. It preserves the entire
consequence relation, not just equations or closed
theorems. Consequently the established full-iPRA
interpretations and equiconsistency transfer unchanged.

This result is different from retaining only progress
introduction and strengthening closure elimination.
Here the original elimination rule is retained, and
the case and progress introductions coexist. It is also
different from allowing arbitrary compound progress,
which would supply stronger induction.

A complete-checker comparison must still count the
certificate choice, its equation-root restriction,
substitution, context handling, computation, and the
public input/output interface. The mathematical theorem
alone is not a byte-count or minimality theorem.
The inference-node bound is distinct from source-size
measurements and from serialized certificate size.
The proof is a written universal proof; finite regression
tests do not constitute a Lean/Rocq formalization of it.


\section{Arithmetic strength and both ordinary interpretation composites}
\label{sec:ordinary-bridge}

We spell out the arithmetic bridge to avoid identifying standard numerical
agreement with bi-interpretability. First elaborate every fixed iterator
into an ordinary function symbol. At stage zero the symbols are $0,S$.
Given a fixed earlier-stage term $h(u,\vec p)$, introduce
$i_h(n,a,\vec p)$ with the two equations
\[
i_h(0,a,\vec p)=a,\qquad
i_h(Sn,a,\vec p)=h(i_h(n,a,\vec p),\vec p).
\]
Only finite descriptions occur; the domain contains no code for a
universal evaluator. Formal parameters are distinct when a description
is declared, while actual arguments can coincide.

Let $T_I$ be ordinary intuitionistic first-order logic with actual
equality and genuine falsity, these staged equations, contextual
induction for literal equations, and
\[
\forall x\,(x=0\vee\exists y\,x=Sy),
\qquad (0=S0)\to\bot.
\]
The second sentence is the falsity bridge. It must be stated: positive
numerical equations by themselves admit a singleton structure. Native
explosion from $0=S0$ and a translation of ordinary falsity to that
equation do not authorize silently dropping the bridge from $T_I$.

The native-to-ordinary translation is
\[
(s=t)^\circ=(s=t),\qquad
(\mathsf D_{x,y}B)^\circ=
\exists x\bigl(B^\circ(x,0)\to\forall y B^\circ(x,y)\bigr).
\]
A clause $\langle A;y.B\rangle$ becomes $A^\circ\to\forall yB^\circ$.
For the auxiliary closure used by the compact checkers,
$(\mathsf Q_yB)^\circ=B^\circ(0)\to\forall yB^\circ$.
The interpretation domain is all elements, defined by $x=x$;
equality is actual equality.

Case DI is sound because numerical coverage combines its retained base
with its successor proof. Progress DI is sound by literal-equation
induction under a temporarily assumed base. Use, DE and transport are
ordinary logical inferences under the translations. Conversely the
derived logical operations, equation induction, cases and explosion in
Section~\ref{sec:choice-introduction} recover the rules and axioms of
$T_I$. An auxiliary closure is interderivable with
$\mathsf D_{v,y}B(y)$ for unused $v$: in one direction use case DI with
the closure and its base; in the other eliminate the paired formula
and return its supplied closure. Thus formulas and context clauses have
the stated meanings in both presentations.

\subsection{The arithmetic dependency order}

The scalar update has no recursion-counter argument. We do not conceal
one in an encoding. The inherited internal proofs, reproduced in the
appendices, establish the required arithmetic in this order:
\begin{enumerate}
\item Using equation induction and numerical cases, construct a fixed
scalar predecessor term and prove $P0=0$ and $P(Sn)=n$.
Existential coverage then gives $x=0\vee x=S(Px)$: in the successor
case $x=Sy$, apply $P$ to obtain $Px=y$ and substitute.
\item Construct truncated subtraction and a numerical characteristic
of actual equality. Prove its reflection on the entire original domain.
\item Compile each fixed quantifier-free formula into its characteristic
equation and recover every contextual quantifier-free induction instance.
\item With that induction available, construct pairing and projections
by a bounded square-root search and compile every fixed general
primitive-recursive description into a scalar iterator term.
\end{enumerate}
No pairing or faithful equality detector is assumed to obtain
predecessor. No appeal to standard-model correctness replaces an
internal equality. Every induction fixes its passive parameters;
diagonal substitution occurs only into a previously proved identity.

For a target recursion
\[
f(0,\vec p)=g(\vec p),\qquad
f(Sn,\vec p)=h(n,f(n,\vec p),\vec p),
\]
let $\pi,L,H$ be the proved pairing and projections. The compiler uses
\[
V(n,\vec p)=I\bigl(n,\pi(0,g^*(\vec p));
u.\pi(SL(u),h^*(L(u),H(u),\vec p))\bigr),
\qquad f^*(n,\vec p)=H(V(n,\vec p)).
\]
Equation induction proves $L(V(n,\vec p))=n$; the two projection
identities then prove the target defining equations. The construction
is external recursion on each finite PR description, not an internal
quantification over arbitrary functions.

\begin{theorem}[Whole-domain bi-interpretability and equiconsistency]
The ordinary realization $T_I$ and full first-order iPRA are
bi-interpretable. In both theories $x=y$ defines the isomorphism from
the original domain to its double interpretation. The native Choice-DI
calculus is therefore a faithful proof presentation of a theory
bi-interpretable and equiconsistent with iPRA.
\end{theorem}
\begin{proof}
Interpret an iterator in iPRA by primitive recursion that ignores its
counter argument in the update. iPRA proves its equations, literal
equation induction, coverage and separation. In the reverse direction
use the scalar compiler. The predecessor and equality-reflection proofs
supply the successor axioms and every quantifier-free induction instance.
All ordinary logical rules are preserved.

For two iterators with the same seed and provably equal updates,
equation induction on the equality of their values proves uniqueness at
every counter, with all parameters fixed. The corresponding argument
applies to ordinary primitive recursion. Induct externally on a finite
description: its double translation satisfies its original defining
equations, so internal uniqueness identifies the two functions. This
also proves coherence when substitution changes the staged syntactic
description of an update.

The domain and equality are unchanged; hence the identity graph is
internally total, bijective and preserves every interpreted operation.
The displayed ordinary meanings of the logical macros and the paired
binder give both formula composites by external induction on finite
formula syntax. Clauses are similarly equivalent to their conditional
universal meanings. The explicit falsity bridge preserves contradiction
in both directions. Effective translations of finite proofs give
equiconsistency, without a claim that either theory proves its own
consistency.
\end{proof}


% Adapted with attribution from the pinned unparenthesized one-hole archive.
\section{Internal scalar predecessor bootstrap}
\label{sec:hybrid-bootstrap}

This scalar bootstrap is inherited from the supplied scalar development~\cite{erasure} and is reproduced to keep the scalar checker self-contained. It is not a new contribution of this paper. The proof uses a factor map fixing odd states and a saturated-iteration construction of predecessor.

The available arithmetic terms are $0$, $S$, and
\[
 I(n,a;u.h),\qquad
 I(0,a;u.h)=a,\qquad
 I(Sn,a;u.h)=h[I(n,a;u.h)/u].
\]
The update binds only $u$.  All other variables in it are fixed numerical
parameters.  We use equality reasoning, the already reconstructed
zero/successor cases, and induction on a single literal equation.
Every induction below keeps its other parameters fixed.  Predecessor,
successor injectivity, pairing, and induction on arbitrary
quantifier-free formulas are not assumptions of the argument.

\subsection{Elementary terms and two finite cycles}

Write $1=S0$ and use larger numerals as successor abbreviations.  Define
\[
 C(n,a,b)=I(n,a;u.b),\qquad
 Z(n)=C(n,1,0),\qquad
 a+b=I(b,a;u.Su).
\]
The update variable of $C$ is absent from its body.  Consequently
\[
 C(0,a,b)=a,\qquad C(Sn,a,b)=b.
\]
Numerical cases show that $Z(n)$ is $0$ or $1$.
The usual commutative-monoid equations for addition follow by literal
equation induction.  For example, $0+b=b$ and $Sa+b=S(a+b)$ are
proved by induction on $b$; these yield commutativity by a further
equation induction, and associativity follows by induction on the last
summand.  We use these equations to reorder finite sums.  The notation
$2a$ means $a+a$ and $3a$ means $a+a+a$; no multiplication operation
is needed for this bootstrap.

Put
\[
 N(n)=I(n,0;u.Z(u)),\qquad
 j(u)=Z(u)+2N(u),\qquad
 M(n)=I(n,0;u.j(u)).
\]
The range of $N$ is contained in $\{0,1\}$: its zero value is $0$,
and its successor value is a value of $Z$.  This uses numerical cases
on its counter, not induction on a disjunction.  Hence
\begin{equation}
 N(n)+Z(N(n))=1,\qquad
 Z(Z(N(n)))=N(n),\qquad N(n+2)=N(n).
 \label{eq:hybrid-parity}
\end{equation}
Here the last equation follows from two successor computations and
the preceding equation.  Literal-equation inductions give
\begin{equation}
 N(2n)=0,\qquad N(2n+1)=1,\qquad N(3n)=N(n).
 \label{eq:hybrid-parity-scaling}
\end{equation}
For the first equation the successor step adds two to the argument of
$N$; the second follows by one computation.  In the third, increasing
$n$ adds three to the left argument and therefore applies $Z$ once to
both sides of the induction equation.

Finite computations now give
\[
 j(0)=1,\qquad j(1)=2,\qquad j(2)=0.
\]
The proof below needs just these three values and
$M(Sn)=j(M(n))$; it does not assume a division algorithm or a theorem
characterizing $M$ as remainder modulo three.

\subsection{A ramp with a global scaling identity}

Define
\[
 G(n)=I(n,0;u.u+M(u)+1).
\]
The central invariant is
\begin{equation}
 M(G(n))=N(n).
 \label{eq:hybrid-ramp-invariant}
\end{equation}
Its zero case is immediate.  For its successor step, retain the
single induction equation and split the Boolean value $N(n)$.
If it is $0$, the next $G$ state is $G(n)+1$; its $M$ value is
$j(0)=1$, equal to $N(Sn)$.  If it is $1$, the next $G$ state
is $G(n)+2$; two computations of $M$ change its value from $1$
to $2$ and then to $0$, again equal to $N(Sn)$.  This proves
the invariant by literal-equation induction.

It follows that
\[
 G(Sn)=G(n)+N(n)+1,\qquad G(n+2)=G(n)+3.
\]
Equation induction therefore yields
\begin{equation}
 G(2n)=3n.
 \label{eq:hybrid-ramp-double}
\end{equation}
A separate literal-equation induction proves
\begin{equation}
 G(3n)=3G(n)+N(n).
 \label{eq:hybrid-ramp-triple}
\end{equation}
For completeness, its successor calculation is
\begin{align*}
 G(3Sn)
 &=G(3n)+N(3n)+4
   =3G(n)+2N(n)+4,\\
 3G(Sn)+N(Sn)
 &=3G(n)+3N(n)+3+Z(N(n)).
\end{align*}
The expressions agree in both Boolean cases for $N(n)$.
The base is $0=0$.  In particular,
\eqref{eq:hybrid-ramp-triple} was proved internally; it was not
inferred from an intended floor-function interpretation of $G$.

Set
\[
 H(u)=C(N(u),G(u),u).
\]
The crucial choice is to retain $u$ when its parity is $1$.
The preceding equations give
\begin{equation}
 H(2u)=3u,\qquad H(1)=1,\qquad H(3u)=3H(u).
 \label{eq:hybrid-factor-map}
\end{equation}
To check the last identity, use $N(3u)=N(u)$ and split the Boolean
$N(u)$.  At $0$, both sides are $3G(u)$ by
\eqref{eq:hybrid-ramp-triple}; at $1$, both are $3u$.
This argument does not require that an arbitrary number has already
been represented as $2v$ or $2v+1$.

\subsection{A comparator from scalar iteration}

For every fixed update $h$, equation induction on $n$, keeping $m,a$
and all parameters fixed, proves the semigroup identity
\[
 I(m+n,a;h)=I(n,I(m,a;h);h).
\]
The base is the zero computation and the step applies $h$ to the
induction equation. This identity will also be used in the later
pairing construction.


For a fixed update $h$, equation induction on $y$, with its initializer
$a$ and all parameters of $h$ fixed, proves the rotation identity
\begin{equation}
 I(Sy,a;h)=I(y,h(a);h).
 \label{eq:hybrid-rotation}
\end{equation}
Both sides have base $h(a)$, and their successor equations apply the
same $h$ to the previous equality.  Here and below $I(y,a;h)$ is
an abbreviation for one fixed finite scalar term, not quantification
over functions.

By \eqref{eq:hybrid-factor-map}, another equation induction gives
\begin{equation}
 I(y,3a;H)=3I(y,a;H).
 \label{eq:hybrid-iterator-scaling}
\end{equation}
Its base is $3a=3a$.  In the step the fixed-parameter induction
equation rewrites the new left side as
\[
 H(3I(y,a;H))=3H(I(y,a;H)).
\]
The initializer parameter stays $a$ throughout this induction.
Equation induction also gives $I(y,1;H)=1$, since $H(1)=1$.

Define
\[
 E(x)=I(x,1;u.2u),\qquad
 L(x,y)=N\bigl(I(y,E(x);H)\bigr).
\]
Then the following entire recursion specification is provable:
\begin{equation}
 L(0,y)=1,\qquad L(Sx,0)=0,\qquad L(Sx,Sy)=L(x,y).
 \label{eq:hybrid-comparator}
\end{equation}
The first equation uses $E(0)=1$ and the fixed-point identity for
$H$.  The second uses $E(Sx)=2E(x)$ and $N(2a)=0$.
For the third,
\begin{align*}
 L(Sx,Sy)
 &=N\bigl(I(Sy,2E(x);H)\bigr)\\
 &=N\bigl(I(y,H(2E(x));H)\bigr)\\
 &=N\bigl(I(y,3E(x);H)\bigr)\\
 &=N\bigl(3I(y,E(x);H)\bigr)\\
 &=N\bigl(I(y,E(x);H)\bigr)=L(x,y).
\end{align*}
The successive steps use rotation, $H(2a)=3a$, iterator scaling,
and $N(3a)=N(a)$.  Thus $L$ has the comparison equations at all
arguments internally.  Its standard-model meaning as the indicator
of $x\le y$ is motivation, not a premise of the following argument.

\subsection{Saturated iteration gives predecessor}

We now need only the three equations
\eqref{eq:hybrid-comparator}.  Define
\[
 B_m(u)=C(L(Su,m),u,Su),\qquad
 W(t,m)=I(t,0;u.B_m(u)).
\]
The subscript in $B_m$ records a fixed parameter; $B_m$ is an open
term abbreviation, not a function variable.

Numerical cases on $e$ prove the selector equation
\[
 C(e,Sa,Sb)=S C(e,a,b).
\]
Equivalently one may use equation induction on $e$, whose induction
hypothesis is unused.  Therefore
\begin{align}
 B_{Sm}(Su)
 &=C(L(SSu,Sm),Su,SSu)\notag\\
 &=C(L(Su,m),Su,SSu)=S B_m(u),
 \label{eq:hybrid-conjugacy}\\
 B_{Sm}(0)&=C(L(S0,Sm),0,S0)=C(L(0,m),0,S0)=S0.
 \label{eq:hybrid-conjugacy-base}
\end{align}

Literal-equation induction on $t$, with $m$ fixed, proves
\begin{equation}
 W(St,Sm)=S W(t,m).
 \label{eq:hybrid-saturated-shift}
\end{equation}
The base is \eqref{eq:hybrid-conjugacy-base}.  The successor step is
\begin{align*}
 W(SSt,Sm)
 &=B_{Sm}(W(St,Sm))\\
 &=B_{Sm}(S W(t,m))\\
 &=S B_m(W(t,m))=S W(St,m).
\end{align*}
Only the second line uses the current induction hypothesis, at the
same fixed $m$.

Two subsequent equation inductions give
\begin{equation}
 L(Sn,n)=0,\qquad W(n,n)=n.
 \label{eq:hybrid-diagonals}
\end{equation}
The first uses the comparator recurrence.  The second uses
\eqref{eq:hybrid-saturated-shift} in its successor step.  The
substitution $t=m=n$ is made into that already proved universally
parameterized equation; no parameter of an active induction hypothesis
is changed.  Consequently
\begin{equation}
 W(Sn,n)=B_n(n)=C(L(Sn,n),n,Sn)=n.
 \label{eq:hybrid-extra-step}
\end{equation}

Finally define the finite scalar term
\begin{equation}
 P(m)=I\bigl(m,0;u.C(L(SSu,m),u,Su)\bigr).
 \label{eq:hybrid-predecessor-term}
\end{equation}
Its zero equation is immediate.  To justify its successor equation
without assuming a separate extensionality rule for iterator bodies,
put
\[
 V(t,m)=I\bigl(t,0;u.C(L(SSu,Sm),u,Su)\bigr).
\]
The comparator recurrence identifies this update pointwise with
$B_m(u)$.  Equation induction on $t$, with $m$ fixed, gives
$V(t,m)=W(t,m)$: their base values are $0$, and the successor step
uses the pointwise update equation and then the induction equation.
Now \eqref{eq:hybrid-extra-step} yields
\[
 \boxed{\qquad P(0)=0,\qquad P(Sm)=V(Sm,m)=W(Sm,m)=m.\qquad}
\]

\subsection{Dependency boundary}

All terms in this section are built from $0,S,I$ by finite substitution.
All inductions have a single equation as their invariant.  Numerical
cases and equality transport are the already reconstructed logical
resources.  Successor injectivity may now be obtained by applying $P$
to an assumed equation $Sx=Sy$.  Equality tests, reflection of
quantifier-free formulas, pairing, and the recovery of general
primitive recursion can then be developed in the established order.
They were not used in order to manufacture predecessor.

This argument is a written internal derivation in the stated theory.
The finite implementation checks do not constitute a proof-assistant
formalization of this entire metatheorem.


% Standalone section.  Requires only amsmath and amssymb in the parent document.
% No theorem environments, macros, or bibliography keys are assumed.
% All arithmetic claims in this section are proved here.

\section{From literal-equation induction to full intuitionistic PRA}
\label{sec:eic-completion}\label{app:arithmetic}

This section supplies an unconditional completion theorem. It is adapted
with attribution from the current forward-left draft~\cite{forward}. In the
scalar application used here, every function occurring in the proof is
already a scalar term: predecessor is supplied by
Appendix~\ref{sec:hybrid-bootstrap}, the selector is
$C(n,a,b)=I(n,a;u.b)$, addition is $x+y=I(y,x;u.Su)$, and
truncated subtraction is $d(x,y)=I(y,x;u.Pu)$. All other functions below
are finite compositions of these. Thus this proof does not presuppose
that arbitrary primitive recursion has already been compiled. The QFI
conclusion first applies to the scalar language; the general PR compiler
comes afterward.  Literal-equation
induction, together with the primitive-recursive defining equations,
predecessor coverage, and separation of $0$ and $1$, proves induction for every
quantifier-free formula.  In particular, the proof does not assume a faithful
numerical equality test, stability or decidability of equality, successor
uniqueness, or any commutation law for an iterator.  We construct the equality
test and prove its reflection property for arbitrary elements of the
numerical domain.

\subsection{The exact hypotheses}

Work in ordinary intuitionistic first-order logic with equality and genuine
falsity $\bot$.  The language has one numerical sort and a symbol for each
fixed, finite primitive-recursive description.  Distinct descriptions are
not identified merely because they compute the same standard function.
Let $\mathsf{PRE}$ contain the defining equations for zero, successor,
projections, composition, and primitive recursion.  In particular, choose
descriptions $P$ and $C$ with equations
\begin{align}
 P(0)&=0, & P(Sx)&=x, \label{eq:eic-predecessor}\\
 C(0,a,b)&=a, & C(Sx,a,b)&=b. \label{eq:eic-selector}
\end{align}
Write $1=S0$.  In addition to $\mathsf{PRE}$, assume
\begin{align}
 \mathsf{Cov}:&\qquad
       \forall x\,\bigl(x=0\lor x=S(Px)\bigr),
       \label{eq:eic-coverage}\\
 \mathsf{Sep}:&\qquad 0=1\longrightarrow\bot,
       \label{eq:eic-separation}
\end{align}
and induction for each \emph{literal equation} $E(n,\vec p)$:
\begin{equation}
 \forall\vec p\,
 \left[
   \left(E(0,\vec p)\land
       \forall n\bigl(E(n,\vec p)\longrightarrow E(Sn,\vec p)\bigr)
   \right)
   \longrightarrow\forall n\,E(n,\vec p)
 \right].
 \label{eq:eic-induction}
\end{equation}
Here the schematic expression $E$ is a single equality of numerical
terms; the displayed universal quantifiers range only over numbers.
Every description and every instance of the schema is finite external
syntax.

The axiom formulation \eqref{eq:eic-induction} and the usual contextual
induction rule are interchangeable.  For example, if
\[
 \Gamma\vdash E(0,\vec p),
 \qquad
 \Gamma,E(z,\vec p)\vdash E(Sz,\vec p),
\]
where $z$ is absent from $\Gamma$ and the passive parameters, discharge the
second premise's equation and universally introduce $z$.  The axiom then
gives $\Gamma\vdash\forall n E(n,\vec p)$, which can be instantiated at any
target term.  Conversely, assume the antecedent of the axiom and apply the
contextual rule with fresh variables.  This uses the ordinary intuitionistic
logical rules and permits arbitrary, including quantified, formulas in
$\Gamma$.

\subsection{Constructor separation and two elementary tests}

Define the sign description and the countdown description by
\begin{equation}
 \operatorname{sg}(u)=C(u,0,1),
 \qquad
 K(u)=C(u,1,Pu).
 \label{eq:eic-sign-countdown}
\end{equation}
Their defining equations imply
\begin{align}
 \operatorname{sg}(0)&=0,
 &\operatorname{sg}(Sx)&=1,\label{eq:eic-sign-computation}\\
 K(0)&=1,
 &K(Sx)&=x.\label{eq:eic-countdown-computation}
\end{align}

\paragraph{Constructor laws.}
Successor is injective: applying $P$ to $Sx=Sy$ gives $x=y$ by
\eqref{eq:eic-predecessor}.  Successor never equals zero.  Indeed,
$Sx=0$ gives $1=0$ by applying $\operatorname{sg}$ and using
\eqref{eq:eic-sign-computation}; symmetry and
\eqref{eq:eic-separation} give $\bot$.  Consequently $1\ne0$, where
$A\ne B$ abbreviates $A=B\to\bot$.  None of these arguments uses
induction.

\paragraph{The sign test.}
For every $u$,
\begin{align}
 &\operatorname{sg}(u)=0\lor\operatorname{sg}(u)=1,
       \label{eq:eic-sign-range}\\
 &\operatorname{sg}(u)=0\quad\longleftrightarrow\quad u=0,
       \label{eq:eic-sign-zero}\\
 &\operatorname{sg}(u)=1\quad\longleftrightarrow\quad u\ne0.
       \label{eq:eic-sign-one}
\end{align}
To prove the range assertion, split \eqref{eq:eic-coverage}.  The branch
$u=0$ gives sign $0$, and the branch $u=S(Pu)$ gives sign $1$.
The implication $u=0\to\operatorname{sg}(u)=0$ is computation.
For its converse, the successor branch of coverage would make
$\operatorname{sg}(u)=0$ imply $1=0$, a contradiction; in that branch
explosion gives $u=0$ as well.  Similarly, if
$\operatorname{sg}(u)=1$ and $u=0$, computation gives $0=1$.
Conversely, $u\ne0$ excludes the zero branch of coverage, and the
successor branch computes the sign as $1$.

Coverage and \eqref{eq:eic-selector} also give, for all $u,v$,
\begin{equation}
 C(u,v,v)=v.
 \label{eq:eic-selector-idempotence}
\end{equation}
Both coverage branches compute to $v$, so this identity again needs no
induction.

\paragraph{The countdown test and the notation $W$.}
For every $u$,
\begin{equation}
 K(u)=0\quad\longleftrightarrow\quad u=1.
 \label{eq:eic-countdown-zero}
\end{equation}
If $u=1=S0$, \eqref{eq:eic-countdown-computation} gives $K(u)=0$.
Conversely, assume $K(u)=0$ and split coverage.  If $u=0$, then
$K(u)=1$, contradicting separation, and explosion gives $u=1$.
If $u=S(Pu)$, computation gives $K(u)=Pu$; hence $Pu=0$ and
$u=S0=1$.

If a presentation uses a description named $W$ with defining equations
\[
 W(0)=1,\qquad W(Sx)=x,
\]
then it proves $K(u)=W(u)$ for every $u$ already from coverage.
In the zero branch both sides are $1$; in the branch $u=S(Pu)$ both
sides are $Pu$.  Thus $K$ and $W$ are provably equal on the whole domain,
even if they are distinct descriptions.  This equality does not appeal to
standard-model agreement or to uniqueness of an arbitrary recurrence.

\subsection{Four elementary uses of equation induction}

Choose descriptions for addition and truncated subtraction, with recursion
on the second argument:
\begin{align}
 x+0&=x, &x+Sy&=S(x+y),\label{eq:eic-addition}\\
 d(x,0)&=x, &d(x,Sy)&=P(d(x,y)).\label{eq:eic-subtraction}
\end{align}
Using these equations and \eqref{eq:eic-induction}, we now derive
\begin{align}
 0+y&=y,\label{eq:eic-left-zero}\\
 (Sx)+y&=S(x+y),\label{eq:eic-left-successor}\\
 d(Sx,Sy)&=d(x,y),\label{eq:eic-cancel-successors}\\
 d(Sa,a)&=1.\label{eq:eic-diagonal}
\end{align}
Each induction below has a single equation as its induction formula.

For \eqref{eq:eic-left-zero}, induct on $y$.  The base is $0+0=0$.
The step is $0+Sy=S(0+y)=Sy$.

For \eqref{eq:eic-left-successor}, hold $x$ fixed and induct on $y$.
At zero both sides are $Sx$.  At a successor the induction hypothesis
gives
\[
 (Sx)+Sy=S\bigl((Sx)+y\bigr)
       =S\bigl(S(x+y)\bigr)=S(x+Sy).
\]
Combining \eqref{eq:eic-left-successor} with the defining equation for
addition gives the useful rearrangement
\begin{equation}
 (Sy)+r=y+Sr.
 \label{eq:eic-shift-successor}
\end{equation}
Indeed, both sides are $S(y+r)$.

For \eqref{eq:eic-cancel-successors}, hold $x$ fixed and induct on $y$.
The base computes as
\[
 d(Sx,S0)=P(d(Sx,0))=P(Sx)=x=d(x,0).
\]
Applying $P$ to the induction equation gives the step:
\[
 d(Sx,S(Sy))=P(d(Sx,Sy))=P(d(x,y))=d(x,Sy).
\]

Finally, induct on $a$ for \eqref{eq:eic-diagonal}.  The base is
$d(S0,0)=S0=1$.  At a successor,
\[
 d(S(Sa),Sa)=d(Sa,a)=1,
\]
using \eqref{eq:eic-cancel-successors} and the induction hypothesis.

\subsection{An equation invariant that reflects actual equality}

The remaining arithmetic ingredient is the following equation, valid for
all $x,y$:
\begin{equation}
 C\bigl(\operatorname{sg}(d(x,y)),\ x,\ y+d(x,y)\bigr)=x.
 \label{eq:eic-residual-invariant}
\end{equation}
Its meaning is that a nonzero subtraction residual reconstructs the
original argument by addition.  The selector expresses this assertion as
one literal equation, so equation induction applies directly.

Hold $x$ fixed and induct on $y$.  At $y=0$, subtraction gives residual
$x$, and \eqref{eq:eic-left-zero} gives $0+x=x$.  The desired equation
is consequently $C(\operatorname{sg}(x),x,x)=x$, which is
\eqref{eq:eic-selector-idempotence}.

For the step, assume \eqref{eq:eic-residual-invariant} at $y$ and
abbreviate
\[
 r=d(x,y),\qquad r'=Pr=d(x,Sy).
\]
These are local term abbreviations, not additional function symbols.
Split coverage for $r'$.

If $r'=0$, then $\operatorname{sg}(r')=0$, and the successor instance
of \eqref{eq:eic-residual-invariant} reduces immediately to $x=x$.

In the other branch, $r'=S(Pr')$, so constructor separation gives
$r'\ne0$.  Split coverage for $r$.  Its zero branch would imply
$r'=P0=0$, a contradiction, from which the desired step equation follows.
In its successor branch,
\[
 r=S(Pr)=Sr',\qquad \operatorname{sg}(r)=1.
\]
The induction equation therefore computes to $y+r=x$.  Meanwhile,
\eqref{eq:eic-shift-successor} gives
\[
 (Sy)+r'=y+Sr'=y+r=x.
\]
Since $r'$ is a successor, $\operatorname{sg}(r')=1$, so this is
exactly the branch selected in the successor instance of
\eqref{eq:eic-residual-invariant}.  This completes the equation induction.
All case distinctions came from coverage; no induction on an implication
or a disjunction was used.

\paragraph{Reflection at residual $1$.}
We have, for arbitrary $a,b$,
\begin{equation}
 d(Sa,b)=1\quad\longleftrightarrow\quad a=b.
 \label{eq:eic-residual-reflection}
\end{equation}
For the forward direction, substitute $x=Sa$, $y=b$ into
\eqref{eq:eic-residual-invariant}.  If $d(Sa,b)=1$, the sign is $1$
and the selector returns its second branch, giving
\[
 b+1=Sa.
\]
The defining equation for addition gives $b+1=S(b+0)=Sb$.  Thus
$Sb=Sa$, and applying $P$ gives $b=a$, hence $a=b$.
Conversely, $a=b$ and \eqref{eq:eic-diagonal} give
$d(Sa,b)=d(Sa,a)=1$.

Define the equality characteristic by the fixed description
\begin{equation}
 e(a,b)=\operatorname{sg}\bigl(K(d(Sa,b))\bigr).
 \label{eq:eic-equality-characteristic}
\end{equation}
The sign, countdown, and residual reflection lemmas give
\begin{align}
 e(a,b)=0&\quad\longleftrightarrow\quad a=b,
       \label{eq:eic-equality-reflection}\\
 e(a,b)=0&\ \lor\ e(a,b)=1.
       \label{eq:eic-equality-range}
\end{align}
In detail, $e(a,b)=0$ is equivalent to $K(d(Sa,b))=0$ by
\eqref{eq:eic-sign-zero}, which is equivalent to $d(Sa,b)=1$ by
\eqref{eq:eic-countdown-zero}, and then to $a=b$ by
\eqref{eq:eic-residual-reflection}.  The range assertion is
\eqref{eq:eic-sign-range} applied to $K(d(Sa,b))$.

The reflection property concerns the logic's actual equality throughout
the domain.  It was not assumed as part of a numerical comparison symbol,
and its proof did not restrict $a$ or $b$ to standard numerals.
It also yields decidable equality.  In the $e(a,b)=0$ branch,
\eqref{eq:eic-equality-reflection} gives $a=b$.  In the $e(a,b)=1$
branch, an assumption $a=b$ would give $e(a,b)=0$ and therefore $0=1$,
contradicting separation.  Hence
\begin{equation}
 \forall a\,\forall b\,\bigl(a=b\lor(a=b\to\bot)\bigr)
 \label{eq:eic-equality-decidable}
\end{equation}
is a theorem, not an extra hypothesis.

\subsection{Compilation of all quantifier-free formulas}

For each fixed quantifier-free formula $A(\vec p)$, construct a term
$b_A(\vec p)$ by external recursion on its finite syntax.  Take equality,
$\bot$, $\land$, $\lor$, and $\to$ as the primitive formula constructors;
negation is implication into $\bot$.  Set
\begin{align}
 b_{s=t}&=e(s,t), & b_{\bot}&=1,\label{eq:eic-boolean-atoms}\\
 b_{A\land D}&=C(b_A,b_D,1),\label{eq:eic-boolean-and}\\
 b_{A\lor D}&=C(b_A,0,b_D),\label{eq:eic-boolean-or}\\
 b_{A\to D}&=C(b_A,b_D,0).\label{eq:eic-boolean-implies}
\end{align}
Every term here is a finite composition of fixed primitive-recursive
descriptions; if desired, projections and composition give it a single
description symbol.  No internal quantification over formulas or programs
is involved.

An external induction on the syntax of $A$ proves, within the arithmetic
theory,
\begin{equation}
 \bigl(b_A=0\lor b_A=1\bigr)
 \quad\text{and}\quad
 \bigl(b_A=0\longleftrightarrow A\bigr).
 \label{eq:eic-boolean-correctness}
\end{equation}
Here are the full connective checks.  In every case $b_A=1$ implies
$\neg A$: if $A$ held, the already established implication
$A\to b_A=0$ would contradict separation.

\paragraph{Atomic formulas and falsity.}
For an equation, both claims are
\eqref{eq:eic-equality-reflection}--\eqref{eq:eic-equality-range}.
For falsity the range assertion is $1=0\lor1=1$, with its right disjunct
proved by reflexivity.  Also $1=0\to\bot$ follows from separation and
symmetry, while $\bot\to1=0$ follows by explosion.  Thus $1=0$ is
equivalent to $\bot$, as required.

\paragraph{Conjunction.}
Split $b_A=0\lor b_A=1$.  In the zero branch, $A$ holds and
$b_{A\land D}=b_D$.  Its range follows from the range of $b_D$,
and its being zero is equivalent to $D$, hence to $A\land D$ under
the known $A$.  In the one branch, $\neg A$ holds and
$b_{A\land D}=1$.  Its range is immediate.  If it were zero,
separation would give $\bot$, hence $A\land D$; conversely an
assumption $A\land D$ contradicts $\neg A$, and explosion gives
$b_{A\land D}=0$.

\paragraph{Disjunction.}
In the branch $b_A=0$, the formula $A$ holds and
$b_{A\lor D}=0$.  The range is immediate, and both $b_{A\lor D}=0$
and $A\lor D$ hold, proving their equivalence.  In the branch
$b_A=1$, we have $\neg A$ and $b_{A\lor D}=b_D$.  Its range follows
from the range of $b_D$.  If its value is zero, then $D$ and hence
$A\lor D$.  Conversely, from $A\lor D$, the $A$ alternative
contradicts $\neg A$ and yields $b_D=0$ by explosion, while the
$D$ alternative yields $b_D=0$ directly.  This proves the required
equivalence in both alternatives.

\paragraph{Implication.}
In the branch $b_A=0$, the formula $A$ holds and
$b_{A\to D}=b_D$.  Its range follows from the range of $b_D$.
If its value is zero, then $D$, which gives $A\to D$.  Conversely,
$A\to D$ and the known $A$ give $D$, hence $b_D=0$.
In the branch $b_A=1$, we have $\neg A$ and
$b_{A\to D}=0$.  Its range is immediate.  To prove $A\to D$,
assume $A$; this contradicts $\neg A$, and explosion gives $D$.
Thus both $b_{A\to D}=0$ and $A\to D$ hold in this branch, and
their equivalence follows.  This argument uses the proved two-value
range of the finite characteristic.  It does not use excluded middle
for arbitrary formulas or change the surrounding intuitionistic logic.

This completes the external syntactic induction proving
\eqref{eq:eic-boolean-correctness} for every quantifier-free formula.

\subsection{Recovery of induction and the full theory}

Let $A(n,\vec p)$ be any fixed quantifier-free formula.  Suppose an
arbitrary retained context $\Gamma$ proves its base and its induction
step:
\[
 \Gamma\vdash A(0,\vec p),
 \qquad
 \Gamma,A(z,\vec p)\vdash A(Sz,\vec p),
\]
with the usual freshness condition on $z$.  By
\eqref{eq:eic-boolean-correctness}, the base gives
$\Gamma\vdash b_A(0,\vec p)=0$.  Under the additional assumption
$b_A(z,\vec p)=0$, the same equivalence gives $A(z,\vec p)$,
the given step gives $A(Sz,\vec p)$, and the equivalence gives
$b_A(Sz,\vec p)=0$.  Apply literal-equation induction to
$b_A(n,\vec p)=0$ and then apply its reflection implication.
The result is $\Gamma\vdash A(t,\vec p)$ for every target term $t$.
Thus the theory proves the full contextual quantifier-free induction
rule, including in arbitrary quantified contexts.

\paragraph{Completion theorem.}
In the common language of fixed primitive-recursive descriptions, the
theory
\[
 \mathsf{PRE}+\mathsf{Cov}+\mathsf{Sep}+\mathsf{EI}
\]
has exactly the same sequents as ordinary first-order intuitionistic
primitive recursive arithmetic with all quantifier-free induction
instances.  The forward simulation follows from the recovered
quantifier-free induction, successor injectivity, and successor separation;
all primitive-recursive defining equations and all logical rules are
shared.  For the reverse simulation, literal-equation induction is an
instance of quantifier-free induction, and successor separation gives
$\mathsf{Sep}$.  Quantifier-free induction proves $\mathsf{Cov}$:
at zero choose its left disjunct, and at a successor use
$Sx=S(P(Sx))$, which follows from $P(Sx)=x$, for its right disjunct.

Both interpretations therefore use the full domain, defined by $x=x$,
and preserve every function symbol, equality, logical connective, and
quantifier.  Both composites have the identity isomorphism, with graph
$x=y$.  The theories are consequently identity-domain bi-interpretable
and equiconsistent.  This completion theorem is unconditional once the
literal-equation induction rule has been obtained; it has no dependency
on a classical successor-uniqueness reduction.


\section{Pairing, primitive recursion, and internal numerical composites}
\label{sec:iterator-coding}
This appendix is adapted with attribution from the progress-closure
arithmetic development~\cite{progress}. Its prerequisites have now been
proved: scalar predecessor, arithmetic reflection, and full quantifier-free
induction. Pairing is a derived proof tool; it is not part of any checker.

\subsection{Pairing from a bounded square-root search}

Once QFI is available, a familiar square-based injection gives a
particularly direct compiler. Write $a\le b$ for $d(a,b)=0$.
The usual linear-order and semiring identities follow by fixed
quantifier-free inductions from the successor, addition, and monus
equations. Multiplication itself is the scalar term
\[
                        a\cdot b=I(b,0;u.u+a).
\]
Only primitive-recursive arithmetic is involved in the following
bounded root search; there is no minimization over an unbounded domain.

For a fixed parameter $n$, put
\[
 U_n(v)=C(d((Sv)^2,n),Sv,v),\qquad
 R_n(i)=I(i,0;v.U_n(v)),\qquad q(n)=R_n(n).
\]
The update increments $v$ exactly when $(v+1)^2\le n$, and
otherwise retains it. QFI in $i$, with $n$ fixed, proves
\[
 R_n(i)\le i,\qquad R_n(i)^2\le n,\qquad
 R_n(i)=i\ \vee\ n<(S(R_n(i)))^2.              \tag{Root invariant}
\]
The base has $R_n(0)=0$. If the update increments, its test
contradicts the strict-upper-bound alternative in the third invariant.
Thus the old state equals $i$, and the new state equals $Si$;
its square bound is the increment test. If the update does not increment, the
strict upper square bound holds and the state remains fixed.
An already established strict upper square bound forces every
subsequent update to remain fixed. These are exactly the cases
needed for the three displayed invariants.

At $i=n$, either the strict upper square bound already holds or
$R_n(n)=n$, in which case $n<(n+1)^2$ supplies it. Thus
\[
                        q(n)^2\le n<(q(n)+1)^2.           \tag{Root}
\]
Monotonicity of squares shows uniqueness of a number satisfying
Root. All formulas just used are quantifier free in the scalar
language, with fixed $n$ during the induction.

Now define
\[
 \pi(a,b)=(a+b)^2+a,\qquad
 L(n)=d(n,q(n)^2),\qquad H(n)=d(q(n),L(n)).
\]
For $k=a+b$, one has $k^2\le k^2+a<(k+1)^2$ since $a\le k$.
Root uniqueness gives $q(\pi(a,b))=a+b$, and therefore
\[
                    L(\pi(a,b))=a,\qquad H(\pi(a,b))=b.
\]
Surjectivity of $\pi$ is unnecessary. The original numerical domain
is never replaced by the range of this coding; pairing is used only
to store the finite state of a fixed primitive recursion. This
standard square pairing is a convenient reconstruction here, not
a claimed new arithmetic basis.

\subsection{Compiling each primitive recursion}

Suppose the fixed target description is
\[
 f(0,\vec p)=g(\vec p),\qquad
 f(Sn,\vec p)=h(n,f(n,\vec p),\vec p).
\]
Assume by external structural induction that the earlier $g,h$
already have scalar translations $g^*,h^*$. Define
\[
 W(u,\vec p)=\pi(SL(u),h^*(L(u),H(u),\vec p)),
\]
\[
 V(n,\vec p)=I(n,\pi(0,g^*(\vec p));u.W(u,\vec p)),
 \qquad f^*(n,\vec p)=H(V(n,\vec p)).
\]
The update has only its state variable bound, with every $p_i$
fixed. Equation induction proves $L(V(n,\vec p))=n$.
The two projection equations and scalar computation then prove
both defining equations for $f^*$ internally. Initial functions
and composition translate directly. This completes the translation
of every fixed finite PR description without adding a universal
evaluator.

\subsection{Why both numerical round trips are internal}

Two fixed scalar iterators with the same seed and provably equal
updates are internally equal at all counters: use literal-equation
induction comparing their values. Their base agrees, and their
successor equations plus congruence prove the step. The same
uniqueness argument works for ordinary primitive recursion.

For an ordinary PR description, induct externally on its finite
construction. The double translation satisfies its original
defining equations, so internal equation induction identifies it
with the original operation. For a native scalar description,
earlier translated updates agree by the external induction
hypothesis; scalar uniqueness identifies the double-translated
iterator. These proofs hold at arbitrary internal counters.

Thus the domain relation $x=y$ is provably a total bijection
preserving every fixed numerical operation in both composites.
The formula constructions and round trips in
Section~\ref{sec:choice-introduction} give the corresponding formula
composites. Both interpretations preserve falsity, taking native
$0=1$ to an ordinary contradiction and ordinary $\bot$ to native
$0=1$. Finite proof translations therefore establish equiconsistency.


\clearpage\section{The complete strict canonical checker}

The following normalized module is independently executable. Its full cost appears in Section~\ref{sec:implementation}.


\begin{lstlisting}[language=Python]

import json
import sys

def fail():
    raise ValueError('invalid certificate')

def rec(n, a, h):
    return [n, a, 0 if n == 0 else h]

def successor(a):
    return [0, a, 0]

def term(t, n):
    if type(t) is int:
        return 0 <= t <= n
    if type(t) is not list:
        return False
    return len(t) == 3 and term(t[0], n) and term(t[1], n) and term(t[2], n + 1) and (t[0] != 0 or t[2] == 0)

def formula(a, n):
    if type(a) is not list or not a:
        return False
    if a[0] == '=':
        return len(a) == 3 and term(a[1], n) and term(a[2], n)
    return len(a) == 2 and a[0] in ('Q', 'D') and formula(a[1], n + (1 if a[0] == 'Q' else 2))

def shift_term(t, amount=1, depth=0):
    if type(t) is int:
        return t + amount if t > depth else t
    return rec(shift_term(t[0], amount, depth), shift_term(t[1], amount, depth), shift_term(t[2], amount, depth + 1))

def shift_formula(a, amount=1, depth=0):
    if a[0] == '=':
        return ['=', shift_term(a[1], amount, depth), shift_term(a[2], amount, depth)]
    return [a[0], shift_formula(a[1], amount, depth + (1 if a[0] == 'Q' else 2))]

def subst_term(t, images, depth=0):
    if type(t) is int:
        if t <= depth:
            return t
        index = t - depth
        return shift_term(images[index - 1], depth) if index <= len(images) else t - len(images)
    return rec(subst_term(t[0], images, depth), subst_term(t[1], images, depth), subst_term(t[2], images, depth + 1))

def subst_formula(a, images, depth=0):
    if a[0] == '=':
        return ['=', subst_term(a[1], images, depth), subst_term(a[2], images, depth)]
    return [a[0], subst_formula(a[1], images, depth + (1 if a[0] == 'Q' else 2))]

def reverse(a):
    return ['=', a[2], a[1]] if a[0] == '=' else a

def derive(proof, context, scope):
    if type(proof) is not list or not proof or type(proof[0]) is not str:
        fail()
    tag = proof[0]
    if tag == 'h' and len(proof) == 2:
        index = proof[1]
        if type(index) is int and 0 <= index < len(context):
            return context[index]
        fail()
    if tag == 'r' and len(proof) == 4:
        n, a, h = proof[1:]
        if not (term(n, scope) and term(a, scope) and term(h, scope + 1)):
            fail()
        return ['=', rec(successor(n), a, h), subst_term(h, [rec(n, a, h)])]
    if tag == 'x' and len(proof) == 4:
        template, major, minor = proof[1:]
        if not formula(template, scope + 1) or template[0] != '=':
            fail()
        equation = derive(major, context, scope)
        if equation[0] != '=' or equation[1] == 0:
            fail()
        if derive(minor, context, scope) != subst_formula(template, [equation[1]]):
            fail()
        return reverse(subst_formula(template, [equation[2]]))
    if tag == 'q' and len(proof) == 4:
        witness, major, base = proof[1:]
        if not term(witness, scope):
            fail()
        closure = derive(major, context, scope)
        if closure[0] != 'Q':
            fail()
        if derive(base, context, scope) != subst_formula(closure[1], [0]):
            fail()
        return reverse(subst_formula(closure[1], [witness]))
    if tag == 'e' and len(proof) == 4:
        goal, major, branch = proof[1:]
        if not formula(goal, scope):
            fail()
        paired = derive(major, context, scope)
        if paired[0] != 'D':
            fail()
        extended = [shift_formula(a) for a in context] + [['Q', paired[1]]]
        if derive(branch, extended, scope + 1) != shift_formula(goal):
            fail()
        return goal
    if tag == 'i' and len(proof) == 4:
        goal, witness, branch = proof[1:]
        if not formula(goal, scope) or goal[0] != 'D' or (not term(witness, scope)):
            fail()
        body = shift_formula(goal[1], 1, 2)
        lifted_witness = shift_term(witness)
        eta = 1 if goal[1][0] == '=' and witness == 0 else 0
        hypothesis = subst_formula(body, [eta, lifted_witness])
        target = subst_formula(body, [successor(1), lifted_witness])
        extended = [shift_formula(a) for a in context] + [hypothesis]
        if derive(branch, extended, scope + 1) != target:
            fail()
        return goal
    fail()

def verify(data):
    if type(data) is not dict or set(data) != {'goal', 'proof'}:
        fail()
    if not formula(data['goal'], 0):
        fail()
    if derive(data['proof'], [], 0) != data['goal']:
        fail()
    return True
if __name__ == '__main__':
    try:
        verify(json.load(sys.stdin))
    except Exception:
        print('REJECT')
        sys.exit(1)
    print('ACCEPT')


\end{lstlisting}

\section{The separate compact scalar byte profile}

This implements the existing guarded calculus, with its documented public-input contract. It is not the canonical-$R$ program.


\begin{lstlisting}[language=Python]

import json

def c(a, g=[0], k=0):
    i = lambda a, t=c, d=0: [i(b, t, d + 5 - len(a)) for b in a] if a * 0 != 0 else a if k <= a < d else d % 2 and i(t, d=3)
    if i(a) in g:
        return
    m, s, u, t = q()
    a = (m, a) == (s and [s, u] or [t] * 3, i(t, u)) and i(t, s) or a
    h = c(m, g, k) or m * 3 or (u and [[s and [s], u, t], i(t, [s, u, t], 3) if s else u]) or (a := i(*a, t, 2))
    k -= 1
    c(i(a, [k]), g + [i(h, (m or ~len(a) % 2) and k)], k)
q = json.load(open(0)).pop
c(q())
print('ACCEPT')


\end{lstlisting}

\clearpage
\begin{thebibliography}{9}
\bibitem{repo} Vladimir Reshetnikov and ProveIt contributors.
\emph{Concise presentations of intuitionistic primitive recursive arithmetic}.
Snapshot \texttt{8eb71c1e570e07ef45ead686628e77f18ef34798}.
\href{https://github.com/VladimirReshetnikov/ProveIt/tree/8eb71c1e570e07ef45ead686628e77f18ef34798/Logic/PrimitiveRecursiveArithmetic/ConciseIPRA}{Pinned repository collection}.
\bibitem{shared} ProveIt contributors. Successor-seeded and shared-continuation research lineage,
\file{Research/successor_seed_completion/theory/prelogic.tex}
and \file{Research/successor_seed_completion/theory/bootstrap.tex}, in the snapshot above.
\bibitem{erasure} ProveIt contributors. Scalar predecessor bootstrap,
\file{Research/certified_zero/reference/scalar_predecessor.tex},
with its inherited attribution, in the snapshot above.
\bibitem{forward} ProveIt contributors. Equality reflection and quantifier-free induction,
\file{Research/certified_zero/reference/equality_reflection.tex},
with its inherited attribution, in the snapshot above.
\bibitem{progress} ProveIt contributors. Primitive-recursion compiler,
\file{Research/certified_zero/reference/primitive_recursion.tex},
with its inherited progress-closure attribution, in the snapshot above.
\bibitem{finite} ProveIt contributors. Finite-control and relational obstructions,
\file{Research/local_scopes/theory/finite_profiles.tex} and
\file{relational_obstruction.tex}, in the snapshot above.
\bibitem{kernel} ProveIt contributors. Native kernel comparisons,
\file{Lean/ConciseIPRA/NativeMinimal.lean} and
\file{NativeReversal.lean}, and the ordinary scalar bridge,
\file{NativeBiInterpretation.lean}, in the snapshot above.
\end{thebibliography}
\end{document}
